% Part I: Foundations
% This file is \include'd from main.tex

\chapter{Introduction}
\label{ch:intro}

\section{The fundamental problem}
\label{sec:intro:problem}

Consider the following four artifacts, drawn from four different domains.

\begin{enumerate}[label=(\roman*)]
\item A Python function implementing binary search uses the loop guard
  \texttt{while low < high} instead of \texttt{while low <= high}.
  Every branch of the loop body is locally correct; the error is
  visible only at the \emph{interface} between the guard and the body.

\item The English sentence ``The horse raced past the barn fell'' is
  locally grammatical in every sub-phrase---``the horse raced past the
  barn'' parses as a relative clause, and ``fell'' is a well-formed
  verb---but the global parse fails because the two fragments cannot
  be composed into a single syntactic tree without additional
  structure.

\item A sonnet whose three quatrains are metrically and imagistically
  perfect but whose volta (the turn between the octave and the sestet)
  introduces a non-sequitur.  Each stanza, read in isolation, is
  accomplished verse; the poem as a whole is incoherent.

\item A mathematical proof in which each lemma is impeccably
  established but the overall argument is circular: Lemma~3 depends on
  Lemma~1, Lemma~1 depends on Lemma~5, and Lemma~5 depends on
  Lemma~3.  Every \emph{local} step is valid; the \emph{global}
  derivation is unsound.
\end{enumerate}

In every case the pattern is the same: \emph{local correctness does
not imply global correctness}.  Locally verified parts fail to compose
into a globally verified whole.  The bug, the ambiguity, the
incoherence, and the circularity are all invisible when one inspects
the parts in isolation; they emerge only at the \emph{overlaps}
between parts.

This observation is the starting point of \textbf{Judgment Geometry}
(JuGeo).  We develop a mathematical framework in which:

\begin{itemize}
\item The artifact under analysis is modeled as a \emph{site}---a
  category of regions equipped with a Grothendieck topology specifying
  which families of sub-regions jointly observe everything about a
  larger region.

\item Properties of the artifact are modeled as \emph{presheaves} (or
  \emph{sheaves}) on the site: functors that assign to each region a
  set of local data and to each inclusion a restriction map.

\item The question ``do locally verified properties hold globally?''
  is answered by the \emph{descent condition}: local sections glue
  into a global section if and only if the first \v{C}ech cohomology
  group vanishes, $\check{H}^1(\mathcal{U}, \mathcal{F}) = 0$.

\item When the cohomology does not vanish, the non-trivial classes are
  \emph{obstructions}---first-class mathematical objects that identify
  exactly which overlaps between regions fail to reconcile.

\item Every claim carries a \emph{grade} drawn from a semiring
  $\Grade$, measuring the quality or confidence of the evidence
  supporting it.
\end{itemize}

The result is a \emph{judgment}: a structured datum that records not
only what is claimed and where, but how well the claim is supported,
what evidence underlies it, what obstacles prevent global validity, and
how those obstacles might be resolved.


\section{Why existing frameworks fail}
\label{sec:intro:existing}

Three broad families of tools address the problem of artifact
correctness, and each captures only a fragment of the full picture.

\paragraph{Type theory and proof assistants.}
Systems such as Lean, Coq, and F$^\star$ are founded on dependent type
theory.  They provide absolute certainty about the claims they
verify---if a term type-checks, the proposition it inhabits is
true---but they operate in a binary regime: either the proof
type-checks or it does not.  There is no room for partial evidence,
for claims supported at different confidence levels by different
sources, or for graded quality.  A flawless proof and a dubious
conjecture are equally rejected when they fail to type-check; a
beautiful proof and a tedious case-bash are equally accepted when they
do.  Moreover, the local-to-global composition problem is solved by
the programmer manually: the system checks individual lemmas but does
not automatically determine whether the lemma graph is globally
coherent.

\paragraph{Testing and runtime analysis.}
Test suites, fuzzing, and runtime monitoring sample the behavior of an
artifact on finitely many inputs.  They can detect bugs but cannot
certify their absence.  They have no internal notion of geometric
structure: a test suite does not know that two test cases exercise
overlapping code regions, or that a passing test on the loop body and
a passing test on the loop guard jointly cover the loop only if the
guard--body interface is itself tested.  There is no notion of
\emph{covering} or of \emph{descent}.

\paragraph{Machine learning and LLM-based analysis.}
Statistical models can classify artifacts (``this code is likely
buggy,'' ``this sentence is likely grammatical'') and even propose
repairs, but they provide confidence scores rather than proofs.  The
scores are not compositional: the confidence that function $f$ is
correct and the confidence that function $g$ is correct do not combine
in any principled way into a confidence that the module $\{f, g\}$ is
correct.  There is no algebra of trust, no sheaf-theoretic gluing, and
no cohomological obstruction theory.

\medskip

Judgment Geometry fills the gap by integrating geometric context
(sites and covers), logical evidence (typed proofs and witnesses),
cohomological invariants (\v{C}ech cohomology classes measuring
failure to glue), and quality measures (grades from a semiring) into a
single, unified mathematical framework.


\section{The core insight}
\label{sec:intro:insight}

The central idea can be stated in a single sentence:

\begin{quote}
\itshape
A judgment is a section of a presheaf over a semantic site, carrying
proof-relevant evidence graded by quality, whose global validity is
governed by the vanishing of \v{C}ech cohomology.
\end{quote}

More precisely, we associate to any structured artifact---a program, a
sentence, a poem, a proof---a small category $\mathcal{C}$ whose
objects are the meaningful sub-regions of the artifact and whose
morphisms are the structural relationships (containment, restriction,
transport) between regions.  We equip $\mathcal{C}$ with a
Grothendieck topology $J$ specifying which families of sub-regions
constitute \emph{covers}.  The pair $(\mathcal{C}, J)$ is a
\emph{site}.

Over this site, properties of the artifact become \emph{presheaves}:
contravariant functors $\mathcal{F} : \mathcal{C}^{\op} \to \Set$
assigning to each region $U$ a set $\mathcal{F}(U)$ of local claims,
with restriction maps that specialize claims to sub-regions.  A
\emph{local section} is an element $s \in \mathcal{F}(U)$---a claim
about a particular region.  A \emph{global section} is a compatible
family of local sections: one for each member of a covering family,
agreeing on all overlaps.

The \v{C}ech cohomology group $\check{H}^1(\mathcal{U}, \mathcal{F})$
measures the obstruction to gluing: it vanishes if and only if every
compatible family of local sections extends to a global section.  When
it does not vanish, its non-trivial elements identify precisely the
overlaps at which gluing fails.

Finally, each claim carries a \emph{grade} $T \in \Grade$, where
$\Grade$ is a semiring measuring quality along multiple dimensions.
The grade algebra provides the quantitative layer: not merely ``is
this artifact correct?'' but ``how good is it, along which axes, and
how does quality compose?''


\section{The 4-object approach}
\label{sec:intro:4objects}

The full data of a judgment can be organized into four maximally rich
mathematical objects, each drawn from a well-established branch of
mathematics:

\[
\JJ = (\STop,\; \PStr,\; \CSpec,\; \GEval).
\]

\begin{description}[leftmargin=2.5em]
\item[$\STop$ --- the Semantic Topos.]
  The site $(\mathcal{C}, J)$ together with the category of sheaves
  $\Sh(\mathcal{C}, J)$.  This is where the artifact lives: its
  regions, its covers, and the gluing conditions that govern
  local-to-global transitions.

\item[$\PStr$ --- the Proof Structure.]
  The dependent-type-theoretic datum: the context, the proposition,
  the evidence term, and the typing judgment $\Gamma \vdash e : A$.
  This is the logical content of the claim: what is asserted, and what
  evidence supports it.

\item[$\CSpec$ --- the Cohomological Spectrum.]
  The \v{C}ech cohomology groups $\check{H}^n(\mathcal{U},
  \mathcal{F})$ and the obstruction classes they contain.  This is the
  diagnostic layer: it measures failure to glue, localizes
  inconsistencies, and guides repair.

\item[$\GEval$ --- the Graded Evaluation.]
  The grade semiring $\Grade$ together with the evaluation functor
  $\mu$ that maps derivations to quality measures.  This is the
  quantitative layer: it assigns a multi-dimensional quality score to
  every judgment and every composition of judgments.
\end{description}

The central thesis of this monograph is that these four objects are
not merely a convenient packaging of the data involved in judgment
but are \emph{exactly the right level of abstraction}: each
corresponds to a deep mathematical structure (topos theory,
type theory, cohomology, quantale theory), and together they
capture everything needed for judgment across domains.

Part~I of this monograph (the present Part) develops the mathematical
foundations: sites, presheaves, sheaves, \v{C}ech cohomology, and the
Grade semiring.  It introduces the judgment and its individual
components.  Part~II then develops each of the four objects
$(\STop, \PStr, \CSpec, \GEval)$ in depth, showing how they unify
these foundations into a single coherent framework.


\section{Motivating examples}
\label{sec:intro:examples}

We return to the four artifacts from \S\ref{sec:intro:problem} and
sketch how Judgment Geometry analyzes each.

\begin{example}[A buggy program]
\label{ex:intro:program}
The binary search function with guard \texttt{while low < high}
determines a site with three coordinates: $c_{\mathrm{guard}}$ (the
loop guard), $c_{\mathrm{body}}$ (the loop body), and
$c_{\mathrm{fn}}$ (the function).  The covering family is
$\{c_{\mathrm{guard}} \to c_{\mathrm{fn}},\; c_{\mathrm{body}} \to
c_{\mathrm{fn}}\}$.  The guard's local section claims ``when
$\mathit{low} = \mathit{high}$, the loop exits without checking
$\mathit{arr}[\mathit{low}]$.''  The body's local section assumes
``$\mathit{arr}[\mathit{mid}]$ is always checked.''  These disagree on
the overlap $c_{\mathrm{guard}} \cap c_{\mathrm{body}}$: the
coboundary $\delta^0 \neq 0$, so $\check{H}^1 \neq 0$.  The
obstruction class identifies the guard--body interface, and the repair
frontier suggests changing \texttt{<} to \texttt{<=}.
\end{example}

\begin{example}[An ambiguous sentence]
\label{ex:intro:sentence}
The garden-path sentence ``The horse raced past the barn fell'' determines
a site whose coordinates are the morphemes, phrases, and the full
sentence.  At the phrase level, ``the horse raced past the barn''
parses as a noun phrase with a reduced relative clause, and ``fell'' is
a well-formed verb phrase.  But on the overlap---the interface where
the parsed noun phrase must serve as the subject of ``fell''---the two
local parses are incompatible: the noun phrase has already consumed the
verb ``raced,'' leaving no room for ``fell'' as the main verb.  The
\v{C}ech cohomology does not vanish; the obstruction localizes the
ambiguity to the phrase--sentence boundary.
\end{example}

\begin{example}[A poem with forced rhyme]
\label{ex:intro:poem}
Consider a sonnet in which the first quatrain is metrically perfect
and imagistically vivid, but the second quatrain distorts its syntax
to force a rhyme: ``And through the misty vale the river \emph{gleams}
/ Reflecting back her silver-threaded \emph{beams}.''  Each quatrain,
analyzed as a local section over the stanza coordinate, receives a
high grade for meter and rhyme but a lower grade for naturalness.  The
grade semiring captures this: $T = T_{\text{meter}} \otG
T_{\text{rhyme}} \otG T_{\text{imagery}} \otG T_{\text{naturalness}}$,
and the forced rhyme depresses $T_{\text{naturalness}}$, pulling the
overall grade down through the bottleneck principle.
\end{example}

\begin{example}[A flawed proof]
\label{ex:intro:proof}
A mathematical argument consists of five lemmas, each individually
correct, but with a circular dependency: Lemma~3 uses Lemma~1,
Lemma~1 uses Lemma~5, and Lemma~5 uses Lemma~3.  The site has
coordinates for each lemma and for the overall theorem.  Each local
section (lemma) is valid; the \v{C}ech computation reveals a
non-trivial 1-cocycle supported on the triangle of overlaps
$\{1 \cap 3,\; 3 \cap 5,\; 5 \cap 1\}$.  More precisely, the
dependency cycle creates a failure of the cocycle condition on triple
overlaps, yielding a non-trivial class in $\check{H}^1$.  The
obstruction pinpoints the circular dependency and the repair frontier
suggests breaking the cycle.
\end{example}


\section{Overview of the monograph}
\label{sec:intro:overview}

This monograph is organized into six Parts.

\paragraph{Part~I: Foundations
  (Chapters~\ref{ch:intro}--\ref{ch:grade}).}
The present Part develops the mathematical foundations from scratch.
Chapter~\ref{ch:sites} introduces category theory, Grothendieck
topologies, presheaves, sheaves, and the Yoneda lemma.
Chapter~\ref{ch:judgment} defines the judgment and its components.
Chapter~\ref{ch:cech} develops \v{C}ech cohomology and descent.
Chapter~\ref{ch:grade} introduces the Grade semiring and the
evaluation functor.

\paragraph{Part~II: The Judgment Framework.}
Develops $\STop$, $\PStr$, $\CSpec$, and $\GEval$ as standalone
mathematical structures, explaining the design principles behind
the 4-tuple and establishing the judgment fibration.  Also shows
that the component-wise flattened view
$(c, \varphi, A, E, O, B, T, \Pi)$ loses no information.

\paragraph{Part~III: Compositional Theory of Theories.}
Extends Judgment Geometry to natural language and poetry via a
multi-stratum architecture in which each linguistic level
(phonological, morphological, syntactic, semantic, pragmatic,
discourse, poetic, aesthetic) is a coordinate in a master site.

\paragraph{Part~IV: Formal Proofs.}
Collects the detailed proofs of all major theorems stated in
Parts~I--III.

\paragraph{Part~V: Worked Examples.}
Presents extended worked examples across all four domains: program
verification, natural language analysis, poetry evaluation, and
mathematical reasoning.

\paragraph{Part~VI: Applications and Future Directions.}
Explores applications of Judgment Geometry to automated repair,
synthesis, equivalence checking, incremental analysis, documentation
alignment, and multi-agent verification.

\chapter{Semantic Sites and Presheaves}
\label{ch:sites}

This chapter develops the categorical foundations on which Judgment
Geometry rests.  We begin with the basic language of category theory,
introduce Grothendieck topologies and sites, and then build presheaves
and sheaves over sites.  The Yoneda lemma, the central organizing
principle of the theory, receives a full proof.  We close with
concrete examples of sites drawn from all four target domains.

The treatment is self-contained: no prior knowledge of category theory
is assumed beyond elementary set theory.

%% ─── Section 1: Category Theory Preliminaries ─────────────────────

\section{Categories, functors, and natural transformations}
\label{sec:sites:categories}

\begin{definition}[Category]
\label{def:category}
A \emph{category} $\mathcal{C}$ consists of:
\begin{enumerate}[label=(\roman*)]
\item A collection $\ob(\mathcal{C})$ of \emph{objects}.
\item For each pair of objects $A, B \in \ob(\mathcal{C})$, a set
  $\Hom_{\mathcal{C}}(A, B)$ of \emph{morphisms} (or \emph{arrows})
  from $A$ to $B$.
\item For each triple $A, B, C$, a \emph{composition} map
  \[
    \circ : \Hom_{\mathcal{C}}(B, C) \times \Hom_{\mathcal{C}}(A, B)
    \to \Hom_{\mathcal{C}}(A, C).
  \]
\item For each object $A$, an \emph{identity} morphism
  $\id_A \in \Hom_{\mathcal{C}}(A, A)$.
\end{enumerate}
These data satisfy:
\begin{itemize}
\item \emph{Associativity}: $h \circ (g \circ f) = (h \circ g) \circ f$
  whenever the compositions are defined.
\item \emph{Identity}: $f \circ \id_A = f = \id_B \circ f$ for every
  $f : A \to B$.
\end{itemize}
A category is \emph{small} if $\ob(\mathcal{C})$ is a set (not a
proper class).
\end{definition}

\begin{example}[Basic categories]
\label{ex:basic-categories}
\begin{enumerate}[label=(\alph*)]
\item $\Set$: the category whose objects are sets and whose morphisms are functions.
\item For any partially ordered set $(P, \leq)$: the category with one morphism $a \to b$ whenever $a \leq b$.
\item For any group $G$: the category with one object $\ast$ and $\Hom(\ast, \ast) = G$, composition given by group multiplication.
\end{enumerate}
\end{example}

\begin{definition}[Functor]
\label{def:functor}
Let $\mathcal{C}$ and $\mathcal{D}$ be categories.  A \emph{functor}
$F : \mathcal{C} \to \mathcal{D}$ consists of:
\begin{enumerate}[label=(\roman*)]
\item A map $F : \ob(\mathcal{C}) \to \ob(\mathcal{D})$.
\item For each pair $A, B \in \ob(\mathcal{C})$, a map
  $F_{A,B} : \Hom_{\mathcal{C}}(A, B) \to \Hom_{\mathcal{D}}(F(A), F(B))$.
\end{enumerate}
These maps preserve identities ($F(\id_A) = \id_{F(A)}$) and
composition ($F(g \circ f) = F(g) \circ F(f)$).

A functor $F : \mathcal{C}^{\op} \to \mathcal{D}$, where
$\mathcal{C}^{\op}$ denotes the \emph{opposite category} (same
objects, reversed morphisms), is called a \emph{contravariant functor}
from $\mathcal{C}$ to $\mathcal{D}$.
\end{definition}

\begin{definition}[Natural transformation]
\label{def:nat-trans}
Let $F, G : \mathcal{C} \to \mathcal{D}$ be functors.  A \emph{natural
transformation} $\eta : F \Rightarrow G$ is a family of morphisms
$\{\eta_A : F(A) \to G(A)\}_{A \in \ob(\mathcal{C})}$ such that for
every morphism $f : A \to B$ in $\mathcal{C}$, the diagram
\[
\begin{tikzcd}
F(A) \ar[r, "\eta_A"] \ar[d, "F(f)"'] & G(A) \ar[d, "G(f)"] \\
F(B) \ar[r, "\eta_B"']                 & G(B)
\end{tikzcd}
\]
commutes: $G(f) \circ \eta_A = \eta_B \circ F(f)$.  If each $\eta_A$
is an isomorphism, $\eta$ is a \emph{natural isomorphism}.
\end{definition}


%% ─── Section 2: Sites ─────────────────────────────────────────────

\section{Grothendieck topologies and sites}
\label{sec:sites:grothendieck}

The notion of a Grothendieck topology generalizes the idea of an open
cover in point-set topology to an arbitrary category.  It specifies,
for each object, which families of morphisms should be regarded as
``covers.''

\begin{definition}[Sieve]
\label{def:sieve}
Let $\mathcal{C}$ be a category and $X \in \ob(\mathcal{C})$.  A
\emph{sieve} on $X$ is a collection $S$ of morphisms with codomain
$X$ that is closed under precomposition: if $(f : U \to X) \in S$ and
$g : V \to U$ is any morphism, then $(f \circ g : V \to X) \in S$.

Equivalently, $S$ is a subfunctor of the representable presheaf
$\Hom_{\mathcal{C}}(-, X) : \mathcal{C}^{\op} \to \Set$.
\end{definition}

\begin{definition}[Grothendieck topology]
\label{def:grothendieck-topology}
A \emph{Grothendieck topology} $J$ on a category $\mathcal{C}$ assigns
to each object $X \in \ob(\mathcal{C})$ a collection $J(X)$ of sieves
on $X$, called \emph{covering sieves}, satisfying three axioms:
\begin{enumerate}
\item \textbf{Identity (maximality).}
  The maximal sieve
  $t_X = \{f : U \to X \mid f \text{ is any morphism to } X\}$
  is in $J(X)$.

\item \textbf{Stability (base change).}
  If $S \in J(X)$ and $g : Y \to X$ is any morphism, then the
  \emph{pullback sieve}
  \[
    g^* S = \{h : Z \to Y \mid g \circ h \in S\}
  \]
  is in $J(Y)$.

\item \textbf{Transitivity (local character).}
  If $S \in J(X)$ and $R$ is a sieve on $X$ such that for every
  $(f : U \to X) \in S$ the pullback sieve $f^* R$ is in $J(U)$, then
  $R \in J(X)$.
\end{enumerate}
A pair $(\mathcal{C}, J)$ is called a \emph{site}.
\end{definition}

\begin{remark}
\label{rem:covering-families}
In practice, one often specifies a topology via \emph{covering
families} rather than sieves.  A covering family of $X$ is a set
$\{f_i : U_i \to X\}_{i \in I}$ of morphisms with common codomain
$X$.  The covering family \emph{generates} the sieve
$S = \{g : V \to X \mid g \text{ factors through some } f_i\}$.
When we say a family ``is a cover,'' we mean the sieve it generates is
a covering sieve.
\end{remark}

\begin{example}[The classical site]
\label{ex:classical-site}
Let $X$ be a topological space and $\mathcal{O}(X)$ the poset category
of open sets (morphisms are inclusions $U \hookrightarrow V$ whenever
$U \subseteq V$).  The Grothendieck topology is defined by declaring
$\{U_i \hookrightarrow U\}_{i \in I}$ to be a covering family whenever
$\bigcup_{i} U_i = U$.  The three axioms are immediate from the
properties of open sets.
\end{example}


%% ─── Section 3: Presheaves ────────────────────────────────────────

\section{Presheaves and the free cocompletion}
\label{sec:sites:presheaves}

\begin{definition}[Presheaf]
\label{def:presheaf}
Let $\mathcal{C}$ be a small category.  A \emph{presheaf} on
$\mathcal{C}$ is a functor $F : \mathcal{C}^{\op} \to \Set$.  The
category of presheaves on $\mathcal{C}$, denoted
\[
  \PSh(\mathcal{C}) = [\mathcal{C}^{\op}, \Set],
\]
has presheaves as objects and natural transformations as morphisms.
\end{definition}

For a morphism $f : U \to V$ in $\mathcal{C}$, a presheaf $F$ gives a
\emph{restriction map} $F(f) : F(V) \to F(U)$.  If $s \in F(V)$ is a
``section over $V$,'' then $F(f)(s) \in F(U)$ is its ``restriction to
$U$.''  We often write $s|_U$ for $F(f)(s)$ when the morphism $f$ is
clear from context.

\begin{definition}[Representable presheaf (Yoneda embedding)]
\label{def:yoneda-embedding}
For each object $c \in \ob(\mathcal{C})$, the \emph{representable
presheaf} $\mathrm{y}(c) = \Hom_{\mathcal{C}}(-, c)$ is defined by:
\begin{align*}
  \mathrm{y}(c)(U) &= \Hom_{\mathcal{C}}(U, c), \\
  \mathrm{y}(c)(f : U \to V) &= (- \circ f) :
    \Hom_{\mathcal{C}}(V, c) \to \Hom_{\mathcal{C}}(U, c).
\end{align*}
The assignment $c \mapsto \mathrm{y}(c)$ defines the \emph{Yoneda
embedding} $\mathrm{y} : \mathcal{C} \to \PSh(\mathcal{C})$.
\end{definition}

\begin{theorem}[The Yoneda lemma]
\label{thm:yoneda}
Let $\mathcal{C}$ be a small category, $c \in \ob(\mathcal{C})$, and
$F \in \PSh(\mathcal{C})$.  There is a bijection, natural in both $c$
and $F$:
\[
  \Hom_{\PSh(\mathcal{C})}(\mathrm{y}(c), F) \;\cong\; F(c).
\]
\end{theorem}

\begin{proof}
We construct explicit maps in both directions and verify they are
mutually inverse.

\medskip\noindent\textbf{From left to right.}  Given a natural
transformation $\eta : \mathrm{y}(c) \Rightarrow F$, define
$\Phi(\eta) = \eta_c(\id_c) \in F(c)$.  This is an element of $F(c)$
because $\eta_c : \mathrm{y}(c)(c) = \Hom(c, c) \to F(c)$ and
$\id_c \in \Hom(c, c)$.

\medskip\noindent\textbf{From right to left.}  Given $x \in F(c)$,
define a natural transformation $\Psi(x) : \mathrm{y}(c) \Rightarrow F$
by setting, for each object $U$ and each morphism $f \in \mathrm{y}(c)(U) = \Hom(U, c)$:
\[
  \Psi(x)_U(f) = F(f)(x) \in F(U).
\]
We must check that $\Psi(x)$ is natural.  Let $g : U \to V$ be a
morphism in $\mathcal{C}$.  For any $f \in \Hom(V, c)$:
\begin{align*}
  (\Psi(x)_U \circ \mathrm{y}(c)(g))(f)
  &= \Psi(x)_U(f \circ g) = F(f \circ g)(x), \\
  (F(g) \circ \Psi(x)_V)(f)
  &= F(g)(F(f)(x)) = F(f \circ g)(x),
\end{align*}
where the last equality uses the functoriality of $F$ (since $F$ is
contravariant, $F(g) \circ F(f) = F(f \circ g)$).  Hence the
naturality square commutes.

\medskip\noindent\textbf{Mutual inverses.}
$\Phi(\Psi(x)) = \Psi(x)_c(\id_c) = F(\id_c)(x) = x$ since $F$
preserves identities.

Conversely, for any $\eta : \mathrm{y}(c) \Rightarrow F$ and any
$f \in \Hom(U, c)$:
\[
  \Psi(\Phi(\eta))_U(f) = F(f)(\eta_c(\id_c))
  = \eta_U(\mathrm{y}(c)(f)(\id_c))
  = \eta_U(f),
\]
where the second equality uses naturality of $\eta$ applied to $f : U \to c$:
\[
  F(f) \circ \eta_c = \eta_U \circ \mathrm{y}(c)(f).
\]
Evaluating both sides at $\id_c$ gives $F(f)(\eta_c(\id_c)) = \eta_U(f)$.
Hence $\Psi(\Phi(\eta)) = \eta$.
\end{proof}

\begin{corollary}[Full faithfulness of Yoneda]
\label{cor:yoneda-ff}
The Yoneda embedding $\mathrm{y} : \mathcal{C} \to \PSh(\mathcal{C})$
is full and faithful: for all $c, d \in \ob(\mathcal{C})$,
\[
  \Hom_{\mathcal{C}}(c, d) \;\cong\;
  \Hom_{\PSh(\mathcal{C})}(\mathrm{y}(c), \mathrm{y}(d)).
\]
\end{corollary}

\begin{proof}
Apply the Yoneda lemma with $F = \mathrm{y}(d)$:
$\Hom(\mathrm{y}(c), \mathrm{y}(d)) \cong \mathrm{y}(d)(c) = \Hom(c, d)$.
\end{proof}

\begin{remark}[Presheaves as the free cocompletion]
\label{rem:free-cocompletion}
The category $\PSh(\mathcal{C})$ is the free cocompletion of
$\mathcal{C}$: it is the universal category with all small colimits,
equipped with a functor $\mathrm{y} : \mathcal{C} \to \PSh(\mathcal{C})$
such that every functor $F : \mathcal{C} \to \mathcal{D}$ into a
cocomplete category $\mathcal{D}$ extends essentially uniquely to a
colimit-preserving functor $\hat{F} : \PSh(\mathcal{C}) \to \mathcal{D}$.
This universal property justifies calling presheaves ``generalized
objects'' of $\mathcal{C}$.
\end{remark}


%% ─── Section 4: Sheaves ──────────────────────────────────────────

\section{Sheaves and sheafification}
\label{sec:sites:sheaves}

A sheaf is a presheaf satisfying a gluing condition with respect to a
Grothendieck topology.  The gluing condition is the categorical
abstraction of the familiar fact that a function defined on a
topological space can be constructed by specifying it on each open set
of a cover, provided the definitions agree on overlaps.

\begin{definition}[Matching family and amalgamation]
\label{def:matching-family}
Let $(\mathcal{C}, J)$ be a site, $F \in \PSh(\mathcal{C})$, and
$\{f_i : U_i \to X\}_{i \in I}$ a covering family in $J(X)$.  A
\emph{matching family} (or \emph{compatible family}) for the cover is
a collection $\{s_i \in F(U_i)\}_{i \in I}$ such that for every pair
$i, j \in I$ and every pair of morphisms
$g : W \to U_i$, $h : W \to U_j$ with $f_i \circ g = f_j \circ h$,
we have $F(g)(s_i) = F(h)(s_j)$.

An \emph{amalgamation} of such a family is a section $s \in F(X)$ with
$F(f_i)(s) = s_i$ for all $i \in I$.
\end{definition}

\begin{definition}[Sheaf]
\label{def:sheaf}
A presheaf $F \in \PSh(\mathcal{C})$ is a \emph{sheaf} for the
topology $J$ if for every covering family $\{f_i : U_i \to X\}$ in $J(X)$,
every matching family has a unique amalgamation.  Equivalently, $F$
is a sheaf if the diagram
\[
  F(X) \;\xrightarrow{\;\sim\;}\;
  \mathrm{eq}\!\left(
    \prod_{i \in I} F(U_i) \;\rightrightarrows\;
    \prod_{(i,j) \in I \times I} F(U_i \times_X U_j)
  \right)
\]
is an isomorphism for every cover, where the two parallel arrows are
induced by the two projections from the fiber product
$U_i \times_X U_j$ to $U_i$ and to $U_j$.

The full subcategory of $\PSh(\mathcal{C})$ consisting of sheaves is
denoted $\Sh(\mathcal{C}, J)$.
\end{definition}

\begin{proposition}[Sheafification]
\label{prop:sheafification}
The inclusion functor $\iota : \Sh(\mathcal{C}, J) \hookrightarrow
\PSh(\mathcal{C})$ has a left adjoint
\[
  a : \PSh(\mathcal{C}) \to \Sh(\mathcal{C}, J),
\]
called \emph{sheafification} (or the \emph{associated sheaf functor}).
For every presheaf $F$, there is a canonical morphism
$\eta_F : F \to \iota(a(F))$ that is universal among morphisms from
$F$ to sheaves.  The adjunction
$a \dashv \iota : \PSh(\mathcal{C}) \rightleftarrows \Sh(\mathcal{C}, J)$
is exact: $a$ preserves all finite limits.
\end{proposition}

\begin{proof}[Proof sketch]
The sheafification $a(F)$ is constructed in two steps.  First, one
forms the ``separated presheaf'' $F^+$ by quotienting each $F(X)$ by
the equivalence relation that identifies two sections $s, t \in F(X)$
whenever there exists a cover $\{U_i \to X\}$ such that $s|_{U_i} = t|_{U_i}$
for all $i$.  This yields a separated presheaf (injectivity of
amalgamation).  Second, one repeats the construction: $(F^+)^+$ is a
sheaf (existence of amalgamation).  Thus $a(F) = (F^+)^+$.  The
adjunction and exactness are standard; see \cite{johnstone2002} for
the full proof.
\end{proof}

\begin{remark}
\label{rem:topos}
The category $\Sh(\mathcal{C}, J)$ is a \emph{Grothendieck topos}: it
is complete, cocomplete, has a subobject classifier, and satisfies
the Giraud axioms.  In Judgment Geometry, the topos $\Sh(\mathcal{C}, J)$
is the first of the four objects, $\STop = \Sh(\mathcal{C}, J)$, and
provides the geometric universe in which judgments live.
\end{remark}


%% ─── Section 5: Examples of sites ─────────────────────────────────

\section{Examples of semantic sites}
\label{sec:sites:examples}

We now illustrate how the abstract machinery of sites applies to the
four domains motivating Judgment Geometry.

\begin{example}[Program verification]
\label{ex:site:programs}
Given a program $P$, the \emph{semantic site}
$\mathcal{S}(P) = (\mathcal{C}_P, J_P)$ is defined as follows:

\medskip\noindent\textbf{Objects (coordinates).}  The objects of
$\mathcal{C}_P$ are the structural units of $P$:
\begin{center}
\begin{tabular}{ll}
\emph{Coordinate kind} & \emph{Program construct} \\
\hline
Module   & top-level file \\
Class    & \texttt{class} definition \\
Function & \texttt{def} / \texttt{async def} \\
Scope    & lexical scope (loop body, \texttt{with}-block) \\
Branch   & \texttt{if}/\texttt{else}/\texttt{elif}/\texttt{match} case
\end{tabular}
\end{center}

\medskip\noindent\textbf{Morphisms.}  Morphisms encode four
structural relationships:
\begin{description}[nosep,leftmargin=2em]
\item[Restriction:] function $\to$ loop body, class $\to$ method.
\item[Inclusion:] branch $\hookrightarrow$ conditional, method $\hookrightarrow$ class.
\item[Transport:] call site $\to$ callee (equivalence-preserving).
\item[Refinement:] coordinate $\to$ finer decomposition.
\end{description}

\medskip\noindent\textbf{Topology.}  Covering families arise from:
\begin{itemize}[nosep]
\item \emph{Control-flow covers}: the branches of an \texttt{if/else}
  cover the conditional.
\item \emph{Scope covers}: a function's local scopes cover the function.
\item \emph{Module covers}: the top-level definitions cover the module.
\end{itemize}

The topology axioms are verified:
\emph{Identity} holds because every coordinate covers itself.
\emph{Stability} holds because restricting a control-flow decomposition
to a sub-scope yields a valid decomposition of the sub-scope.
\emph{Transitivity} holds because nested decompositions compose: if a
module is covered by its functions, and each function is covered by its
branches, the composition covers the module.
\end{example}

\begin{example}[Natural language semantics]
\label{ex:site:nl}
A sentence $S$ determines a site $\mathcal{S}(S)$ whose objects are
linguistic units at multiple granularities:
\begin{center}
\begin{tabular}{ll}
\emph{Coordinate kind} & \emph{Linguistic unit} \\
\hline
Morpheme  & smallest meaningful unit (\emph{un-}, \emph{-ing}) \\
Word      & lexical item \\
Phrase    & syntactic constituent (NP, VP, PP) \\
Clause    & finite or non-finite clause \\
Sentence  & complete sentence \\
Discourse & multi-sentence passage
\end{tabular}
\end{center}
Morphisms are the compositional structure maps: morpheme $\to$ word
(affixation), word $\to$ phrase (phrase-structure rules), phrase $\to$
clause (predication), clause $\to$ sentence (embedding/coordination),
sentence $\to$ discourse (discourse relations).  A word is covered by
its morphemes; a phrase is covered by its immediate constituents; a
sentence is covered by its clauses.

A presheaf over this site assigns to each level a set of semantic
representations (e.g., $\lambda$-terms in the Montagovian tradition),
and the restriction maps implement compositional semantics: the
meaning of a phrase is determined by---and restricts to---the meanings
of its parts.
\end{example}

\begin{example}[Poetry]
\label{ex:site:poetry}
A poem $P$ determines a site $\mathcal{S}(P)$ with a hierarchy of
prosodic and structural units:
\begin{center}
\begin{tabular}{ll}
\emph{Coordinate kind} & \emph{Poetic unit} \\
\hline
Syllable & individual syllable (stressed/unstressed) \\
Foot     & metrical unit (iamb, trochee, dactyl, \ldots) \\
Line     & a line of verse \\
Couplet/Tercet & grouping of 2 or 3 lines \\
Stanza   & stanza (quatrain, sestet, \ldots) \\
Poem     & the entire poem
\end{tabular}
\end{center}
Morphisms encode prosodic containment: syllable $\to$ foot $\to$ line
$\to$ stanza $\to$ poem.  A line is covered by its feet; a stanza is
covered by its lines.  The presheaf assigns to each unit the relevant
quality data: metrical regularity, rhyme adherence, imagery, emotional
tone.  Restriction maps propagate constraints: the meter of a line is
determined by the meters of its feet, and the rhyme scheme of a stanza
is determined by the end-words of its lines.
\end{example}

\begin{example}[Mathematics]
\label{ex:site:math}
A mathematical theory $\mathcal{T}$ determines a site whose objects
are the components of the theory:
\begin{center}
\begin{tabular}{ll}
\emph{Coordinate kind} & \emph{Mathematical unit} \\
\hline
Definition  & definition of a concept \\
Lemma       & auxiliary result \\
Theorem     & main result \\
Proof step  & individual step in a proof \\
Theory      & the entire theory
\end{tabular}
\end{center}
Morphisms encode logical dependency: a theorem depends on (receives
morphisms from) the lemmas it cites; a lemma depends on the
definitions it uses.  A theorem is covered by the lemmas and
definitions that constitute its proof.  The presheaf assigns to each
coordinate the set of logical conditions (axioms, hypotheses) active
at that coordinate, and restriction maps propagate hypotheses down the
dependency tree.
\end{example}


%% ─── Section 6: Restriction and Extension ──────────────────────────

\section{Restriction and extension along morphisms}
\label{sec:sites:restriction}

A fundamental operation in sheaf theory is the \emph{restriction} of
data along morphisms.  Given a presheaf $F$ on $\mathcal{C}$ and a
morphism $f : U \to V$, the restriction map $F(f) : F(V) \to F(U)$
``pulls back'' data from $V$ to $U$.  In all four domains, restriction
has a concrete interpretation:

\begin{itemize}
\item \textbf{Programs:} restricting a module-level specification to a
  function yields the function's local verification conditions.

\item \textbf{Natural language:} restricting a sentence-level semantic
  representation to a phrase yields the phrase's meaning in the
  sentential context.

\item \textbf{Poetry:} restricting a poem-level quality assessment to a
  stanza yields the stanza's contribution to the poem's quality.

\item \textbf{Mathematics:} restricting a theory-level consistency
  requirement to a lemma yields the lemma's proof obligation.
\end{itemize}

\begin{proposition}[Restriction preserves compatibility]
\label{prop:restriction-compatibility}
Let $\{f_i : U_i \to X\}$ be a covering family and
$\{s_i \in F(U_i)\}$ a matching family.  If $g : Y \to X$ is any
morphism, then the restricted family $\{F(\mathrm{pr}_i)(s_i)\}$, where
$\mathrm{pr}_i : U_i \times_X Y \to U_i$ is the projection, is a matching
family for the pullback cover $\{U_i \times_X Y \to Y\}$.
\end{proposition}

\begin{proof}
The pullback cover $\{U_i \times_X Y \to Y\}_{i \in I}$ is a cover by
the stability axiom.  The compatibility of the restricted family follows
from the functoriality of $F$: if $s_i|_W = s_j|_W$ on the overlap
$W = U_i \times_X U_j$, then restricting further along $W \times_X Y$
preserves this agreement because $F$ respects composition of morphisms.
\end{proof}

\begin{remark}[Extension]
\label{rem:extension}
The dual operation, \emph{extension} (or \emph{pushforward}) along a
morphism, is more subtle: not every local section extends to a larger
region.  The obstruction to extension is precisely the content of
\v{C}ech cohomology, which we develop in Chapter~\ref{ch:cech}.
The distinction between restriction (always possible) and extension
(possible only when cohomological obstructions vanish) is the
mathematical heart of the local-to-global problem.
\end{remark}

\chapter{The Judgment and Its Components}
\label{ch:judgment}

With the language of sites and presheaves in hand, we now define the
central algebraic object of Judgment Geometry: the \emph{judgment}.
A judgment is not merely a logical claim (``$A$ is true'') nor merely
a typed term (``$e : A$'').  It is a structured datum that records
\emph{where} the claim is made, \emph{what} data is being judged,
\emph{what} is claimed about it, \emph{what evidence} supports the
claim, \emph{what obstacles} prevent global validity, \emph{what
cohomological invariants} measure failure to glue, \emph{how good}
the evidence is, and \emph{how the quality assessment was computed}.

We define a judgment as a 4-tuple $\JJ = (\STop, \PStr, \CSpec,
\GEval)$ of structured objects.  Each of these objects has internal
components, and when we flatten the structured objects to their most
visible constituents, we obtain an 8-component view
$(c, \varphi, A, E, O, B, T, \Pi)$ that is useful for concrete
computation.  This chapter introduces the components; Part~II develops
each structured object in full depth.

\section{Definition of the judgment}
\label{sec:judgment:def}

\begin{definition}[Judgment]
\label{def:judgment}
Let $(\mathcal{C}, J)$ be a site and $\mathcal{F}$ a presheaf on
$\mathcal{C}$.  A \emph{judgment} is a 4-tuple
\[
  \JJ = (\STop,\; \PStr,\; \CSpec,\; \GEval)
\]
whose four objects encode the \textbf{Where}, \textbf{What},
\textbf{Why not}, and \textbf{How good} of the claim (developed
fully in Part~II).

Each object has internal components.  When we flatten these to their
most visible constituents, we obtain the \emph{component-wise view}:
\[
  (c,\; \varphi,\; A,\; E,\; O,\; B,\; T,\; \Pi)
\]
where:
\begin{enumerate}[label=(\roman*),itemsep=4pt]
\item $c \in \ob(\mathcal{C})$ is the \textbf{site component}: the
  coordinate (region, locus) in the site to which the judgment
  applies.

\item $\varphi \in \mathcal{F}(c)$ is the \textbf{presheaf component}:
  the data being judged, as a section of the presheaf $\mathcal{F}$
  at coordinate $c$.

\item $A$ is the \textbf{type component}: the specification,
  proposition, or type against which $\varphi$ is judged.  In a
  program verification context, $A$ might be a postcondition; in
  natural language, a grammaticality constraint; in poetry, a metrical
  template; in mathematics, a theorem statement.

\item $E$ is the \textbf{evidence component}: the proof, witness, or
  test result supporting the claim $\varphi : A$.  Evidence is
  \emph{typed}: we write $E : A$ in context $c$ to indicate that $E$
  is a witness for $A$ at coordinate $c$.  Evidence may be a formal
  proof, an SMT solver certificate, a runtime test trace, or an
  expert assertion.

\item $O$ is the \textbf{obstruction component}: an open cover
  $\mathcal{U} = \{U_i \to c\}$ witnessing non-trivial topology.
  The cover $O$ records the decomposition of $c$ into sub-regions
  over which local verification was performed.  If $c$ is atomic
  (admits only the trivial cover), we set $O = \{c \to c\}$.

\item $B \in \check{H}^1(\mathcal{U}, \mathcal{F})$ is the
  \textbf{cohomology class}: the \v{C}ech 1-class measuring failure
  of the local sections to glue over the cover $O$.  If
  $B = 0$, the local verifications compose into a global
  verification.  If $B \neq 0$, the class identifies which overlaps
  fail and how.

\item $T \in \Grade$ is the \textbf{grade}: a quality measure drawn
  from the Grade semiring (\Cref{ch:grade}), recording how well the
  evidence supports the claim.  The grade is \emph{not} a probability;
  it is an element of an algebraic structure that supports
  composition, comparison, and residuation.

\item $\Pi$ is the \textbf{proof trail}: a record of how $T$ was
  computed---which evidence channels were consulted, which rules were
  applied, which compositions were performed.  The proof trail
  provides auditability: any grade can be traced back to its
  evidential sources.
\end{enumerate}

\noindent The components cluster naturally into four pairs---$(c,
\varphi)$, $(A, E)$, $(O, B)$, $(T, \Pi)$---each pair being the
visible face of one of the four structured objects $\STop$, $\PStr$,
$\CSpec$, $\GEval$.
\end{definition}

\begin{remark}[Relation to Martin-L\"of judgments]
\label{rem:mlof}
Martin-L\"of's type theory uses judgments of the form
$\Gamma \vdash a : A$ (three components: context, term, type).
F$^\star$ extends this to $\Gamma \vdash e : M\;a\;\mathit{wp}$ (four
components: context, term, monad, weakest precondition).  The JuGeo
judgment subsumes both: setting $O$ to the trivial cover, $B = 0$,
$T = \gone$, and $\Pi = \emptyset$ recovers a Martin-L\"of judgment
from the remaining triple $(c, \varphi, A)$ together with the
evidence $E$.  The additional components---obstruction, cohomology
class, grade, and proof trail---carry the geometric, diagnostic, and
qualitative information that type theory alone does not track.
\end{remark}


\section{The components in detail}
\label{sec:judgment:components}

\subsection{The site component \texorpdfstring{$c$}{c}}
\label{sec:judgment:site-component}

The coordinate $c$ answers the question: \emph{where does this
judgment apply?}  A judgment without a coordinate is meaningless---it
is a claim about nothing in particular.  By anchoring every judgment
to a coordinate in the site, we gain:

\begin{itemize}
\item \textbf{Locality:} we can restrict a judgment to sub-regions
  and ask whether it holds there.
\item \textbf{Compositionality:} we can compose judgments along shared
  boundaries (overlaps between coordinates).
\item \textbf{Geometric structure:} the covering families of the site
  organize the local-to-global analysis.
\end{itemize}

\subsection{The presheaf component \texorpdfstring{$\varphi$}{φ}}
\label{sec:judgment:presheaf-component}

The section $\varphi \in \mathcal{F}(c)$ is the data being judged.
In different domains:

\begin{itemize}
\item \textbf{Programs:} $\varphi$ is the verification condition at
  coordinate $c$---a logical formula expressing what must hold for the
  code at $c$ to satisfy its specification.
\item \textbf{Natural language:} $\varphi$ is the semantic
  representation at $c$---a $\lambda$-term, discourse referent, or
  feature bundle.
\item \textbf{Poetry:} $\varphi$ is the quality profile at $c$---a
  record of metrical stress, rhyme, imagery, and emotional tone.
\item \textbf{Mathematics:} $\varphi$ is the logical content at
  $c$---the statement of a lemma, the hypothesis set of a proof step.
\end{itemize}

The key property of $\varphi$ is that it transforms covariantly under
restriction: if $f : U \to c$ is a morphism (so $U$ is a sub-region of
$c$), then $\varphi|_U = \mathcal{F}(f)(\varphi) \in \mathcal{F}(U)$
is the restriction of $\varphi$ to $U$.

\subsection{The type component \texorpdfstring{$A$}{A}}
\label{sec:judgment:type-component}

The type $A$ specifies \emph{what it means for $\varphi$ to be
correct} at coordinate $c$.  Following the Curry-Howard correspondence,
$A$ can be viewed simultaneously as:
\begin{itemize}
\item A \emph{type} that evidence must inhabit: $E : A$.
\item A \emph{proposition} that evidence must prove.
\item A \emph{specification} that the artifact must satisfy.
\end{itemize}

Types in Judgment Geometry are classified into five \emph{proposition
kinds}:

\begin{description}[leftmargin=2.5em]
\item[Structural:] properties of shape---type signatures, arity,
  return types, syntactic well-formedness.
\item[Behavioral:] properties of execution---postconditions,
  preconditions, loop invariants, termination.
\item[Relational:] properties relating multiple coordinates---equivalence,
  commutativity, idempotence, refinement.
\item[Resource:] properties about resource usage---purity, effect
  tracking, memory safety.
\item[Semantic:] higher-level properties---liveness, confluence,
  coherence, aesthetic quality.
\end{description}

The proposition kind determines which evidence channels are applicable,
which descent strategies are appropriate, and how the grade is computed.

\subsection{The evidence component \texorpdfstring{$E$}{E}}
\label{sec:judgment:evidence}

Evidence is a \emph{multi-channel} datum: a judgment may be supported
by evidence from several sources simultaneously.  Formally, $E$ is a
multiset of evidence items, each tagged with a \emph{channel}:

\begin{description}[leftmargin=2.5em]
\item[Proof:] a formal derivation in a proof calculus.
\item[Solver:] a certificate from an SMT solver (Z3, CVC5, etc.)
  operating within a decidable fragment.
\item[Runtime:] an execution trace on a declared set of inputs.
\item[Oracle:] an assertion from an external agent (human expert, LLM).
\end{description}

The typing judgment $E : A$ in context $c$ asserts that $E$ is
\emph{admissible} evidence for $A$ at $c$.  Different channels produce
evidence at different confidence levels, and these levels are
tracked by the grade $T$.

\subsection{The obstruction component \texorpdfstring{$O$}{O}}
\label{sec:judgment:obstruction}

The cover $O = \{U_i \to c\}_{i \in I}$ records the decomposition of
the coordinate $c$ into sub-regions over which local analysis was
performed.  The cover serves two purposes:

\begin{enumerate}
\item It specifies the \emph{granularity} of the analysis: a fine cover
  (many small patches) gives more detailed information but requires
  checking more overlaps; a coarse cover (few large patches) is cheaper
  but may miss local inconsistencies.

\item It provides the input to the \v{C}ech complex
  (\Cref{ch:cech}): the cohomology class $B$ is computed relative to
  the cover $O$.  Different covers may yield different cohomology
  groups; the \emph{true} cohomology is the colimit over all covers
  (the direct limit over refinements).
\end{enumerate}

\subsection{The cohomology class \texorpdfstring{$B$}{B}}
\label{sec:judgment:cohomology}

The class $B \in \check{H}^1(\mathcal{U}, \mathcal{F})$ is the central
diagnostic datum of the judgment.  It answers the question:
\emph{do the local verifications compose into a global verification?}

\begin{itemize}
\item If $B = 0$: the local sections glue.  The artifact is globally
  correct (relative to the specification $A$ and the evidence $E$).

\item If $B \neq 0$: the local sections fail to glue.  The class $B$
  identifies exactly which overlaps are inconsistent and how.  Its
  \emph{support} $\mathrm{supp}(B) = \{(i, j) \mid B_{ij} \neq 0\}$
  localizes the failure; its structure suggests repairs.
\end{itemize}

The full development of \v{C}ech cohomology and its role in Judgment
Geometry is the subject of \Cref{ch:cech}.

\subsection{The grade \texorpdfstring{$T$}{T}}
\label{sec:judgment:grade}

The grade $T \in \Grade$ is a quality measure.  Unlike a probability
(which satisfies the Kolmogorov axioms and lives in $[0, 1]$ with
additive structure), the grade lives in a \emph{semiring}
$(\Grade, \opG, \otG, \gzero, \gone)$ with two operations:

\begin{itemize}
\item $\opG$ (competition/alternative): takes the better of two grades.
\item $\otG$ (composition/conjunction): combines grades along a chain,
  typically by taking the minimum (bottleneck).
\end{itemize}

The full development of the Grade semiring is the subject of
\Cref{ch:grade}.

\subsection{The proof trail \texorpdfstring{$\Pi$}{Π}}
\label{sec:judgment:trail}

The proof trail $\Pi$ is an append-only record documenting:
\begin{enumerate}[nosep]
\item Which evidence channels were consulted and in what order.
\item Which composition rules were applied (restriction, transport,
  gluing).
\item Which grade computations were performed (and the intermediate
  values).
\item Whether any promotions or demotions of grade occurred, and their
  justifications.
\end{enumerate}

The proof trail provides \emph{auditability}: any stakeholder can
trace a judgment's grade back to its evidential roots.  In security-
and safety-critical applications, the proof trail is the certification
artifact.


\section{Composition of judgments}
\label{sec:judgment:composition}

Judgments compose along morphisms in the site.  This composition is the
algebraic content of the gluing operation.

\begin{definition}[Restriction of judgments]
\label{def:judgment-restriction}
Let $\JJ = (c, \varphi, A, E, O, B, T, \Pi)$ be a judgment and
$f : U \to c$ a morphism in the site.  The \emph{restriction} of
$\JJ$ along $f$ is
\[
  \JJ|_U = (U,\; \varphi|_U,\; A|_U,\; E|_U,\; f^*O,\; f^*B,\;
  T',\; \Pi \cdot [\text{restrict along } f])
\]
where:
\begin{itemize}[nosep]
\item $\varphi|_U = \mathcal{F}(f)(\varphi)$,
\item $A|_U$ is the specification specialized to $U$,
\item $E|_U$ is the evidence restricted to $U$,
\item $f^*O$ and $f^*B$ are the pulled-back cover and cohomology class,
\item $T' \leq T$ (restriction may \emph{attenuate} the grade but never
  increase it),
\item $\Pi \cdot [\text{restrict along } f]$ appends the restriction
  step to the proof trail.
\end{itemize}
\end{definition}

\begin{definition}[Composition of judgments]
\label{def:judgment-composition}
Let $\JJ_1 = (c_1, \varphi_1, A_1, E_1, O_1, B_1, T_1, \Pi_1)$ and
$\JJ_2 = (c_2, \varphi_2, A_2, E_2, O_2, B_2, T_2, \Pi_2)$ be
judgments whose coordinates share an overlap: there exist morphisms
$g_1 : W \to c_1$ and $g_2 : W \to c_2$ with
$\varphi_1|_W = \varphi_2|_W$.  The \emph{composition}
$\JJ_1 \circ \JJ_2$ is the judgment
\[
  \JJ_1 \circ \JJ_2 = (c_1 \cup c_2,\; \varphi,\; A_1 \wedge A_2,\;
  E_1 \sqcup E_2,\; O,\; B,\; T_1 \otG T_2,\;
  \Pi_1 \cdot \Pi_2 \cdot [\text{compose}])
\]
where $\varphi$ is the glued section (existing precisely when $B = 0$
for the two-element cover $\{c_1, c_2\}$), $O$ is the union cover,
and $B$ is the \v{C}ech class of the composed cover.
\end{definition}

\begin{proposition}[Associativity of composition]
\label{prop:composition-assoc}
Judgment composition is associative: $(\JJ_1 \circ \JJ_2) \circ \JJ_3
= \JJ_1 \circ (\JJ_2 \circ \JJ_3)$ whenever both sides are defined.
\end{proposition}

\begin{proof}
Associativity follows from the associativity of section gluing (which
is a consequence of the sheaf condition), the associativity of $\otG$
in the Grade semiring, and the associativity of evidence union.
\end{proof}


\section{The judgment category}
\label{sec:judgment:category}

Judgments over a fixed site $(\mathcal{C}, J)$ and presheaf
$\mathcal{F}$ organize into a category.

\begin{definition}[Category of judgments]
\label{def:judgment-category}
The \emph{judgment category} $\mathbf{Jud}(\mathcal{C}, J, \mathcal{F})$
has:
\begin{itemize}
\item \textbf{Objects:} judgments
  $\JJ = (c, \varphi, A, E, O, B, T, \Pi)$.

\item \textbf{Morphisms:} a morphism $\JJ_1 \to \JJ_2$ (called a
  \emph{refinement}) consists of:
  \begin{enumerate}[nosep,label=(\alph*)]
  \item A morphism $f : c_1 \to c_2$ in $\mathcal{C}$ (the judgment
    moves to a finer or coarser coordinate).
  \item A refinement of evidence: $E_2$ extends or strengthens $E_1$
    relative to $f$.
  \item A grade inequality: $T_2 \geq T_1$ (refinement improves or
    preserves quality).
  \item Compatibility: $\varphi_2|_{c_1} = \varphi_1$ and
    $A_2|_{c_1}$ implies $A_1$.
  \end{enumerate}
\end{itemize}
\end{definition}

\begin{proposition}[Well-definedness]
\label{prop:judgment-cat-welldef}
$\mathbf{Jud}(\mathcal{C}, J, \mathcal{F})$ is a category: composition
of refinements is associative and has identities.
\end{proposition}

\begin{proof}
The identity refinement at $\JJ$ is $(\id_c, \id_E, T = T, \id_\varphi)$.
Composition of refinements $(f, E_2, T_2, \ldots) \circ (g, E_3, T_3, \ldots)
= (f \circ g, E_3, T_3, \ldots)$ is associative because composition in
$\mathcal{C}$ is associative and $\geq$ on $\Grade$ is transitive.
\end{proof}

\begin{remark}[Enrichment]
\label{rem:enrichment}
The judgment category is naturally \emph{enriched} over $\Grade$: the
``hom-set'' between two judgments carries a grade measuring the
quality of the refinement.  This enrichment reflects the fact that some
refinements are better than others---a refinement supported by a formal
proof is better than one supported only by an oracle assertion.
The enriched perspective becomes central in Part~II when we define
the Graded Evaluation object $\GEval$.
\end{remark}


\section{The four structured objects}
\label{sec:judgment:structured}

\begin{remark}[From components to structure]
\label{rem:components-to-structure}
The component-wise view $(c, \varphi, A, E, O, B, T, \Pi)$
is useful for concrete calculations, but its eight entries are
\emph{not independent}.  They cluster into four natural groups:

\begin{enumerate}
\item \textbf{Geometric:} $c$ and $O$ (the site component and the
  cover) belong to the site and its topology.  Together with the
  presheaf $\mathcal{F}$ (which supplies $\varphi$), they constitute the
  \emph{Semantic Topos} $\STop$.

\item \textbf{Logical:} $A$ and $E$ (the type and the evidence)
  constitute a dependent typing judgment $E : A$ in context $c$.
  Together they form the \emph{Proof Structure} $\PStr$.

\item \textbf{Cohomological:} $O$ and $B$ (the cover and the
  cohomology class) together with the full \v{C}ech complex
  constitute the \emph{Cohomological Spectrum} $\CSpec$.

\item \textbf{Evaluative:} $T$ and $\Pi$ (the grade and the proof
  trail) together with the Grade semiring constitute the
  \emph{Graded Evaluation} $\GEval$.
\end{enumerate}

These four objects are developed in full depth in Part~II.  Each is
\emph{maximally rich}: $\STop$ is a Grothendieck topos, $\PStr$
is a comprehension category, $\CSpec$ is a spectral sequence, and
$\GEval$ is a quantale with evaluation functor.  The component-wise
view is the shadow of these four objects; the 4-tuple
$\JJ = (\STop, \PStr, \CSpec, \GEval)$ is the natural, canonical
formulation.
\end{remark}

\chapter{\v{C}ech Cohomology and Descent}
\label{ch:cech}

This chapter develops the cohomological heart of Judgment Geometry.
The \v{C}ech complex provides a computable invariant that detects
whether locally verified properties compose into globally verified
ones.  When they do not, the cohomology classes identify exactly where
and how the composition fails.

We develop the theory from first principles, without assuming prior
knowledge of algebraic topology.

%% ─── Section 1: Open covers ──────────────────────────────────────

\section{Open covers and refinements}
\label{sec:cech:covers}

\begin{definition}[Open cover]
\label{def:open-cover}
Let $(\mathcal{C}, J)$ be a site and $X \in \ob(\mathcal{C})$.  An
\emph{open cover} of $X$ is a covering family
$\mathcal{U} = \{f_i : U_i \to X\}_{i \in I}$ belonging to $J(X)$.
\end{definition}

\begin{definition}[Overlaps]
\label{def:overlaps}
Given a cover $\mathcal{U} = \{U_i \to X\}$, the \emph{double overlap}
(or \emph{pairwise intersection}) of $U_i$ and $U_j$ is the fiber
product
\[
  U_{ij} = U_i \times_X U_j,
\]
equipped with projection morphisms $\mathrm{pr}_1 : U_{ij} \to U_i$
and $\mathrm{pr}_2 : U_{ij} \to U_j$.  The \emph{triple overlap} is
\[
  U_{ijk} = U_i \times_X U_j \times_X U_k,
\]
and higher overlaps are defined analogously.
\end{definition}

\begin{definition}[Refinement of covers]
\label{def:cover-refinement}
A cover $\mathcal{V} = \{V_\alpha \to X\}$ is a \emph{refinement} of
$\mathcal{U} = \{U_i \to X\}$ if there exists a map
$\lambda : \mathcal{V} \to \mathcal{U}$ (on index sets) and morphisms
$\rho_\alpha : V_\alpha \to U_{\lambda(\alpha)}$ compatible with the
structure maps to $X$.  The collection of all covers of $X$, ordered
by refinement, forms a directed system.
\end{definition}


%% ─── Section 2: The Čech complex ──────────────────────────────────

\section{The \v{C}ech complex}
\label{sec:cech:complex}

\begin{definition}[\v{C}ech cochain groups]
\label{def:cech-cochain}
Let $\mathcal{U} = \{U_i \to X\}_{i \in I}$ be a cover and
$\mathcal{F}$ a presheaf on $(\mathcal{C}, J)$.  The \emph{\v{C}ech
cochain groups} are:
\begin{align}
  \Cech^0(\mathcal{U}, \mathcal{F})
    &= \prod_{i \in I} \mathcal{F}(U_i)
    && \text{(local sections on each patch)}, \\
  \Cech^1(\mathcal{U}, \mathcal{F})
    &= \prod_{i < j} \mathcal{F}(U_{ij})
    && \text{(data on double overlaps)}, \\
  \Cech^2(\mathcal{U}, \mathcal{F})
    &= \prod_{i < j < k} \mathcal{F}(U_{ijk})
    && \text{(data on triple overlaps)}.
\end{align}
More generally, $\Cech^n(\mathcal{U}, \mathcal{F}) =
\prod_{i_0 < \cdots < i_n} \mathcal{F}(U_{i_0 \cdots i_n})$.
\end{definition}

The key to the entire theory is the coboundary map, which measures
\emph{discrepancy} between overlapping local sections.

\begin{definition}[Coboundary maps]
\label{def:coboundary}
The \emph{zeroth coboundary map}
$\delta^0 : \Cech^0(\mathcal{U}, \mathcal{F})
\to \Cech^1(\mathcal{U}, \mathcal{F})$
sends a family of local sections $(s_i)_{i \in I}$ to the family of
discrepancies:
\[
  (\delta^0 s)_{ij}
  = s_i|_{U_{ij}} - s_j|_{U_{ij}}
  \in \mathcal{F}(U_{ij}).
\]
Here the ``subtraction'' is the comparison operation in $\mathcal{F}$:
$(\delta^0 s)_{ij}$ measures how much the restrictions of $s_i$ and $s_j$
to their overlap disagree.  If $\mathcal{F}$ takes values in abelian
groups, this is literal subtraction; in general, it is the natural
comparison in the fiber.

The \emph{first coboundary map}
$\delta^1 : \Cech^1(\mathcal{U}, \mathcal{F})
\to \Cech^2(\mathcal{U}, \mathcal{F})$
sends a family $(\alpha_{ij})$ to:
\[
  (\delta^1 \alpha)_{ijk}
  = \alpha_{jk}|_{U_{ijk}} - \alpha_{ik}|_{U_{ijk}} + \alpha_{ij}|_{U_{ijk}}.
\]
\end{definition}

\begin{proposition}[$\delta^1 \circ \delta^0 = 0$]
\label{prop:delta-squared-zero}
The composition $\delta^1 \circ \delta^0 = 0$: for any family
$(s_i) \in \Cech^0$,
\[
  (\delta^1 \delta^0 s)_{ijk} = 0.
\]
\end{proposition}

\begin{proof}
Compute directly:
\begin{align*}
  (\delta^1 \delta^0 s)_{ijk}
  &= (\delta^0 s)_{jk}|_{U_{ijk}}
   - (\delta^0 s)_{ik}|_{U_{ijk}}
   + (\delta^0 s)_{ij}|_{U_{ijk}} \\
  &= (s_j|_{U_{jk}} - s_k|_{U_{jk}})|_{U_{ijk}}
   - (s_i|_{U_{ik}} - s_k|_{U_{ik}})|_{U_{ijk}}
   + (s_i|_{U_{ij}} - s_j|_{U_{ij}})|_{U_{ijk}} \\
  &= s_j|_{U_{ijk}} - s_k|_{U_{ijk}}
   - s_i|_{U_{ijk}} + s_k|_{U_{ijk}}
   + s_i|_{U_{ijk}} - s_j|_{U_{ijk}} \\
  &= 0. \qedhere
\end{align*}
\end{proof}

The sequence $\Cech^0 \xrightarrow{\delta^0} \Cech^1 \xrightarrow{\delta^1} \Cech^2 \to \cdots$
is therefore a \emph{cochain complex}: the image of each coboundary
map is contained in the kernel of the next.

\begin{remark}[Interpretation of the degrees]
\label{rem:cech-degrees}
\begin{itemize}
\item An element of $\Cech^0$ is a \emph{choice of local data}:
  one section on each patch.
\item An element of $\Cech^1$ is \emph{overlap data}: one datum on
  each pairwise intersection.
\item An element of $\Cech^2$ is \emph{higher compatibility data}:
  one datum on each triple intersection.
\end{itemize}
The coboundary $\delta^0$ measures \emph{pairwise discrepancy}; the
coboundary $\delta^1$ measures \emph{failure of the cocycle condition}
on triple overlaps.
\end{remark}


%% ─── Section 3: Čech cohomology ──────────────────────────────────

\section{\v{C}ech cohomology groups}
\label{sec:cech:cohomology}

\begin{definition}[\v{C}ech cohomology]
\label{def:cech-cohomology}
The \emph{\v{C}ech cohomology groups} of $\mathcal{F}$ with respect to
the cover $\mathcal{U}$ are:
\[
  \check{H}^n(\mathcal{U}, \mathcal{F})
  = \frac{\ker(\delta^n : \Cech^n \to \Cech^{n+1})}
         {\operatorname{im}(\delta^{n-1} : \Cech^{n-1} \to \Cech^n)},
\]
with the convention that $\Cech^{-1} = 0$.  The elements of
$\ker \delta^n$ are called \emph{$n$-cocycles}; the elements of
$\operatorname{im} \delta^{n-1}$ are called \emph{$n$-coboundaries}.
\end{definition}

The first three cohomology groups have particularly clear
interpretations.

\subsection{\texorpdfstring{$\check{H}^0$}{H0}: Global sections}
\label{sec:cech:h0}

\begin{proposition}[$H^0$ is the group of global sections]
\label{prop:h0-global-sections}
$\check{H}^0(\mathcal{U}, \mathcal{F}) = \ker \delta^0$
consists of those families of local sections $(s_i)_{i \in I}$
that agree on all pairwise overlaps:
\[
  s_i|_{U_{ij}} = s_j|_{U_{ij}} \quad \text{for all } i, j.
\]
When $\mathcal{F}$ is a sheaf, such a compatible family glues to a
unique global section $s \in \mathcal{F}(X)$, so
$\check{H}^0(\mathcal{U}, \mathcal{F}) \cong \mathcal{F}(X)$.
\end{proposition}

\begin{proof}
The condition $\delta^0(s) = 0$ means $s_i|_{U_{ij}} - s_j|_{U_{ij}} = 0$
for all $i, j$, which is exactly the matching family condition of
\Cref{def:matching-family}.  By the sheaf condition
(\Cref{def:sheaf}), every matching family has a unique amalgamation,
giving $\check{H}^0 \cong \mathcal{F}(X)$.
\end{proof}

\subsection{\texorpdfstring{$\check{H}^1$}{H1}: Obstructions to gluing}
\label{sec:cech:h1}

This is the most important cohomology group for Judgment Geometry.

\begin{proposition}[$H^1$ measures failure to glue]
\label{prop:h1-obstructions}
A cocycle $\alpha \in \ker \delta^1 \subset \Cech^1$ is a family
$(\alpha_{ij})$ satisfying the \emph{cocycle condition} on triple overlaps:
\[
  \alpha_{ij}|_{U_{ijk}} + \alpha_{jk}|_{U_{ijk}} = \alpha_{ik}|_{U_{ijk}}
  \quad \text{for all } i, j, k.
\]
The cocycle $\alpha$ is a \emph{coboundary} if there exist sections
$t_i \in \mathcal{F}(U_i)$ with $\alpha_{ij} = t_i|_{U_{ij}} - t_j|_{U_{ij}}$,
i.e., $\alpha = \delta^0(t)$.

The quotient $\check{H}^1(\mathcal{U}, \mathcal{F})
= \ker \delta^1 / \operatorname{im} \delta^0$ measures the
\emph{essential} obstructions to gluing: two cocycles that differ by a
coboundary represent the same obstruction, because the coboundary
component can be absorbed by modifying the local sections.

If $\check{H}^1 = 0$, every cocycle is a coboundary, and every
compatible family of local sections (after appropriate modification)
glues to a global section.  If $\check{H}^1 \neq 0$, there exist
genuine obstructions to gluing that cannot be removed by modifying
individual local sections.
\end{proposition}

\subsection{\texorpdfstring{$\check{H}^2$}{H2}: Obstructions to modifying gluings}
\label{sec:cech:h2}

\begin{remark}[$H^2$ and structural obstructions]
\label{rem:h2}
The group $\check{H}^2(\mathcal{U}, \mathcal{F})$ measures obstructions
at a higher level: the failure of attempted gluings to be consistent
across triple overlaps.  In Judgment Geometry, $\check{H}^2 \neq 0$
signals a \emph{structural} obstruction---the cover itself is
inadequate, and the decomposition of the artifact must be redesigned.
This corresponds to architectural refactoring (in programs), paragraph
restructuring (in text), formal restructuring (in poetry), or
reorganization of the proof structure (in mathematics).

In practice, $\check{H}^2$ obstructions are rarer than $\check{H}^1$
obstructions but harder to resolve, because they require changing the
cover rather than merely modifying local sections.
\end{remark}


%% ─── Section 4: Descent ──────────────────────────────────────────

\section{Descent: the local-to-global principle}
\label{sec:cech:descent}

The central theorem of Judgment Geometry connects the vanishing of
\v{C}ech cohomology to the existence of global sections.

\begin{definition}[Descent datum]
\label{def:descent-datum}
A \emph{descent datum} for a presheaf $\mathcal{F}$ relative to a
cover $\mathcal{U} = \{U_i \to X\}$ is a 1-cocycle
$\alpha \in \ker \delta^1 \subset \Cech^1(\mathcal{U}, \mathcal{F})$.
The descent datum is \emph{effective} if $\alpha$ is a coboundary:
$\alpha = \delta^0(s)$ for some $(s_i) \in \Cech^0$.
\end{definition}

\begin{theorem}[Fundamental Theorem of Judgment Geometry]
\label{thm:fundamental}
Let $(\mathcal{C}, J)$ be a site, $\mathcal{F}$ a presheaf on
$\mathcal{C}$, and $X \in \ob(\mathcal{C})$.  Let $A$ be a
specification for the artifact at $X$.  Then:

\begin{enumerate}
\item The artifact at $X$ satisfies $A$ if and only if there exists a
  global section $s \in \mathcal{F}(X)$ witnessing $A$.

\item A global section exists if and only if, for some (equivalently,
  every sufficiently fine) cover $\mathcal{U}$ of $X$, the first
  \v{C}ech cohomology vanishes:
  \[
    \check{H}^1(\mathcal{U}, \mathcal{F}) = 0.
  \]

\item When $\check{H}^1(\mathcal{U}, \mathcal{F}) \neq 0$, each
  non-trivial class $[{\alpha}] \in \check{H}^1$ is an
  \emph{obstruction}: it identifies the overlaps at which local
  verifications fail to compose, and its support
  $\mathrm{supp}([\alpha]) = \{(i,j) \mid \alpha_{ij} \neq 0\}$
  localizes the failure.
\end{enumerate}
\end{theorem}

\begin{proof}
(1)~follows from the definition: a global section is a witness.

(2)~$(\Rightarrow)$: If $s \in \mathcal{F}(X)$ is a global section,
restricting to each $U_i$ gives a compatible family:
$s_i = s|_{U_i}$.  Then $\delta^0(s)_{ij} = s_i|_{U_{ij}} - s_j|_{U_{ij}}
= s|_{U_{ij}} - s|_{U_{ij}} = 0$, so every 1-cocycle is the zero
cocycle (since any cocycle $\alpha$ can be written as $\alpha = \alpha - 0
= \alpha - \delta^0(s)$, and if the presheaf is separated, the coboundary
structure forces $[\alpha] = 0$ in $\check{H}^1$).

$(\Leftarrow)$: If $\check{H}^1 = 0$, every 1-cocycle is a coboundary.
Given local sections $s_i \in \mathcal{F}(U_i)$, the discrepancy
$\delta^0(s)$ is a 1-cocycle (by \Cref{prop:delta-squared-zero}),
hence a coboundary.  This means there exist modified local sections
$s'_i$ with $\delta^0(s') = 0$---a compatible family.  For a sheaf
$\mathcal{F}$, this compatible family has a unique amalgamation:
a global section $s \in \mathcal{F}(X)$.

(3)~follows from the definition of support: $\alpha_{ij} \neq 0$
means the restrictions of $s_i$ and $s_j$ to $U_{ij}$ disagree, and
this disagreement cannot be eliminated by modifying individual sections
(since $[\alpha] \neq 0$).
\end{proof}


%% ─── Section 5: The exact sequence ───────────────────────────────

\section{The \v{C}ech exact sequence}
\label{sec:cech:exact}

The \v{C}ech cohomology groups fit into a fundamental exact sequence
that organizes the local-to-global analysis.

\begin{proposition}[The \v{C}ech exact sequence]
\label{prop:cech-exact}
For any cover $\mathcal{U} = \{U_i \to X\}$ and presheaf $\mathcal{F}$,
there is an exact sequence:
\[
  0 \to \check{H}^0(\mathcal{U}, \mathcal{F})
  \xrightarrow{\;\iota\;}
  \prod_{i} \mathcal{F}(U_i)
  \xrightarrow{\;\delta^0\;}
  \prod_{i<j} \mathcal{F}(U_{ij})
  \xrightarrow{\;\delta^1\;}
  \prod_{i<j<k} \mathcal{F}(U_{ijk})
  \to \cdots
\]
The map $\iota$ sends a global section $s$ to the family of
restrictions $(s|_{U_i})$.  Exactness at $\prod \mathcal{F}(U_i)$
states that a family $(s_i)$ is in the image of $\iota$ if and only if
$\delta^0(s) = 0$---precisely the sheaf condition.
\end{proposition}


\section{Long exact sequences from short exact sequences}
\label{sec:cech:les}

\begin{theorem}[Long exact sequence in \v{C}ech cohomology]
\label{thm:cech-les}
Let $0 \to \mathcal{F}' \to \mathcal{F} \to \mathcal{F}'' \to 0$ be a
short exact sequence of presheaves on $(\mathcal{C}, J)$.  Then for
any cover $\mathcal{U}$ there is a long exact sequence:
\[
\begin{tikzcd}[column sep=small]
  0 \ar[r] &
  \check{H}^0(\mathcal{U}, \mathcal{F}') \ar[r] &
  \check{H}^0(\mathcal{U}, \mathcal{F}) \ar[r] &
  \check{H}^0(\mathcal{U}, \mathcal{F}'') \ar[dll, "\delta"',
    out=-10, in=170, looseness=0.7] \\
  & \check{H}^1(\mathcal{U}, \mathcal{F}') \ar[r] &
  \check{H}^1(\mathcal{U}, \mathcal{F}) \ar[r] &
  \check{H}^1(\mathcal{U}, \mathcal{F}'') \ar[dll, "\delta"',
    out=-10, in=170, looseness=0.7] \\
  & \check{H}^2(\mathcal{U}, \mathcal{F}') \ar[r] &
  \cdots &
\end{tikzcd}
\]
The connecting homomorphism $\delta : \check{H}^n(\mathcal{U}, \mathcal{F}'')
\to \check{H}^{n+1}(\mathcal{U}, \mathcal{F}')$ is the obstruction to
lifting: a cohomology class in $\mathcal{F}''$ lifts to $\mathcal{F}$
if and only if its image under $\delta$ vanishes.
\end{theorem}

\begin{proof}[Proof sketch]
The short exact sequence of presheaves induces a short exact sequence
of cochain complexes:
$0 \to \Cech^\bullet(\mathcal{U}, \mathcal{F}') \to
\Cech^\bullet(\mathcal{U}, \mathcal{F}) \to
\Cech^\bullet(\mathcal{U}, \mathcal{F}'') \to 0$.
The long exact sequence in cohomology then follows from the standard
snake lemma argument in homological algebra.  The connecting
homomorphism $\delta$ is constructed by diagram chasing: given a
cocycle $\alpha'' \in \Cech^n(\mathcal{U}, \mathcal{F}'')$, lift it
to $\alpha \in \Cech^n(\mathcal{U}, \mathcal{F})$, apply $\delta^n$
to get an element of $\Cech^{n+1}(\mathcal{U}, \mathcal{F})$, and
verify that this element lies in the image of
$\Cech^{n+1}(\mathcal{U}, \mathcal{F}')$.  The class of this element
is $\delta[\alpha'']$, and it is independent of the choice of lift.
\end{proof}

\begin{remark}[Operational significance]
\label{rem:les-operational}
The long exact sequence has immediate practical consequences.  If
$\mathcal{F}'$ is the presheaf of structural properties (type
signatures, arity) and $\mathcal{F}''$ is the presheaf of behavioral
properties (postconditions, invariants), then $\mathcal{F} = \mathcal{F}'
\oplus \mathcal{F}''$ is the full judgment presheaf.  The exact
sequence tells us:

\begin{enumerate}[nosep]
\item If structural verification succeeds ($\check{H}^1(\mathcal{F}') = 0$)
  but behavioral verification fails ($\check{H}^1(\mathcal{F}'') \neq 0$),
  then full verification fails:
  $\check{H}^1(\mathcal{F}) \neq 0$.

\item The converse need not hold: behavioral obstructions may cancel
  with structural data via the connecting homomorphism.

\item Structural verification can be performed first (it is cheaper),
  and behavioral verification deferred until structural verification
  succeeds.  The exact sequence guarantees this staging is sound.
\end{enumerate}
\end{remark}


%% ─── Section 6: Examples ─────────────────────────────────────────

\section{Examples of descent and obstruction}
\label{sec:cech:examples}

We now present four worked examples, one from each domain, illustrating
how \v{C}ech cohomology detects and localizes failures of global
coherence.

\begin{example}[Programs: circular imports]
\label{ex:cech:programs}
Consider a Python project with three modules $M_1$, $M_2$, $M_3$:
\begin{itemize}[nosep]
\item $M_1$ imports a function from $M_2$.
\item $M_2$ imports a class from $M_3$.
\item $M_3$ imports a constant from $M_1$.
\end{itemize}
Each module type-checks in isolation (the local sections are valid).
The site has coordinates $\{M_1, M_2, M_3, P\}$ where $P$ is the
project, with covering family $\{M_i \to P\}_{i=1}^{3}$.  The
overlaps $M_{ij} = M_i \cap M_j$ are the shared interfaces (imports).

The \v{C}ech 1-cocycle $\alpha$ is defined by:
\begin{align*}
  \alpha_{12} &= \text{``$M_1$ expects function $f$ from $M_2$''}, \\
  \alpha_{23} &= \text{``$M_2$ expects class $C$ from $M_3$''}, \\
  \alpha_{13} &= \text{``$M_3$ expects constant $k$ from $M_1$''}.
\end{align*}
The cocycle condition on the triple overlap $M_{123}$ requires
$\alpha_{12} + \alpha_{23} = \alpha_{13}$, which fails because the
dependency chain $M_1 \to M_2 \to M_3 \to M_1$ is circular.  Hence
$\check{H}^1 \neq 0$, and the obstruction class is supported on the
triangle $\{(1,2), (2,3), (1,3)\}$, identifying the circular import.
\end{example}

\begin{example}[Natural language: garden-path sentences]
\label{ex:cech:nl}
The sentence ``The horse raced past the barn fell'' has a site with
coordinates for each phrase and the full sentence $S$.  The cover
$\{P_1, P_2\} \to S$ decomposes the sentence into ``the horse raced
past the barn'' ($P_1$, parsed as a relative clause modification of
``horse'') and ``fell'' ($P_2$, the main predicate).

At $P_1$: the local section assigns ``horse'' the semantic role of
patient (the horse that was raced).  At $P_2$: the local section
requires a subject.  On the overlap $P_1 \cap P_2$: $P_1$ has
consumed the subject position (``horse'' is the patient of ``raced''),
but $P_2$ needs it as the agent of ``fell.''  The discrepancy
$\delta^0 \neq 0$: the two local parses make incompatible demands on
the subject.  $\check{H}^1 \neq 0$, localizing the ambiguity to the
phrase--sentence boundary.
\end{example}

\begin{example}[Poetry: failed volta]
\label{ex:cech:poetry}
A Petrarchan sonnet has an octave (lines 1--8) and a sestet
(lines 9--14), connected by the \emph{volta} (the ``turn'' at line 9).
The site has coordinates $\{Q_1, Q_2, V, Q_3, C, P\}$: two quatrains,
the volta, a sestet ``quatrain,'' a couplet, and the poem.  The cover
$\{Q_1, Q_2, V, Q_3, C\} \to P$ decomposes the poem.

Suppose each quatrain develops the theme of loss, but the volta
introduces, without transition, the topic of gardening.  The local
section at $Q_2$ (thematic content: loss) and the local section at
$Q_3$ (thematic content: gardening) disagree on the overlap
$Q_2 \cap Q_3 = V$: the volta should mediate the thematic transition,
but instead introduces a non-sequitur.  The coboundary
$\delta^0 \neq 0$, and $\check{H}^1 \neq 0$, with the obstruction
supported at the volta.
\end{example}

\begin{example}[Mathematics: circular proof]
\label{ex:cech:math}
A proof of Theorem $T$ uses Lemmas $L_1, L_3, L_5$, with
dependencies $L_1 \to L_3 \to L_5 \to L_1$.  The site has coordinates
$\{L_1, L_3, L_5, T\}$ with cover $\{L_i \to T\}$.  Each lemma is
individually valid (local sections exist), but the overlaps encode
the dependency structure.

The 1-cocycle $\alpha$ records the logical dependency on each overlap:
$\alpha_{13}$ records that $L_1$ uses $L_3$; $\alpha_{35}$ records
that $L_3$ uses $L_5$; $\alpha_{15}$ records that $L_5$ uses $L_1$.
The cocycle condition demands $\alpha_{13} + \alpha_{35} = \alpha_{15}$
(transitivity of dependency), but this fails because the dependency
is circular rather than transitive.  $\check{H}^1 \neq 0$; the
obstruction identifies the circular dependency and the repair frontier
suggests breaking the cycle at any one edge.
\end{example}


\section{Descent as verification}
\label{sec:cech:descent-verification}

We summarize the connection between descent and verification.

\begin{theorem}[Descent criterion for verification]
\label{thm:descent-verification}
Let $(\mathcal{C}, J)$ be the semantic site of an artifact, and let
$\mathcal{F}$ be the presheaf of verification conditions.  Checking
$\check{H}^1(\mathcal{U}, \mathcal{F}) = 0$ for a cover $\mathcal{U}$
is equivalent to checking that the artifact is globally correct
relative to the specification.

More precisely, the descent procedure is:
\begin{enumerate}[nosep]
\item \textbf{Local verification:} compute local sections
  $s_i \in \mathcal{F}(U_i)$ for each patch $U_i$.
\item \textbf{Discrepancy computation:} compute the coboundary
  $\delta^0(s)$ to measure pairwise disagreement on overlaps.
\item \textbf{Cohomological classification:}
  \begin{itemize}[nosep]
  \item If $\delta^0(s) = 0$: the local sections agree; glue them into
    a global section.  The artifact is verified.
  \item If $[\delta^0(s)] = 0$ in $\check{H}^1$ but $\delta^0(s) \neq 0$:
    the discrepancy is a coboundary; modify local sections to achieve
    agreement, then glue.
  \item If $[\delta^0(s)] \neq 0$ in $\check{H}^1$: the artifact has
    genuine obstructions.  Report the obstruction classes, their
    supports, and repair frontiers.
  \end{itemize}
\end{enumerate}
\end{theorem}

\begin{remark}
The descent procedure is \emph{algorithmic}: given a finite site (as
in all practical applications), the \v{C}ech complex is finite, and
the cohomology groups are computable.  The computational cost is
$O(|I|^2 \cdot T_{\mathrm{restrict}})$, where $|I|$ is the number of
patches and $T_{\mathrm{restrict}}$ is the cost of a single restriction
operation.  For typical artifacts, $|I|$ is small (2--20 patches), so
the quadratic factor is negligible; the dominant cost is the
verification of individual patches.
\end{remark}

\chapter{The Grade Semiring}
\label{ch:grade}

The preceding chapters developed the geometric (sites, covers) and
cohomological (\v{C}ech complex, descent) layers of Judgment Geometry.
These layers answer the qualitative question: \emph{does the artifact
pass or fail?}  This chapter introduces the quantitative layer: the
\emph{Grade semiring}, which answers the question \emph{how good is
the artifact, along which axes, and how does quality compose?}

The Grade semiring is not an ad hoc scoring system bolted onto the
framework.  It is a principled algebraic structure---a quantale---with
two operations (competition and composition), a partial order, and a
residuation operation that internalizes implication.  The semiring
interacts tightly with the site and the cohomology: grades restrict
along morphisms, compose along chains, and attenuate through descent.


%% ─── Section 1: Definition ────────────────────────────────────────

\section{Definition of the Grade semiring}
\label{sec:grade:def}

\begin{definition}[Grade semiring]
\label{def:grade-semiring}
The \emph{Grade semiring} is a tuple
\[
  \Grade = ([0,1],\; \opG,\; \otG,\; \gzero,\; \gone)
\]
where:
\begin{enumerate}[label=(\roman*),itemsep=3pt]
\item The \emph{carrier} is the closed unit interval $[0, 1]$.

\item $\opG : [0,1] \times [0,1] \to [0,1]$ is the
  \textbf{competition} (or \textbf{alternative}) operation:
  $a \opG b = \max(a, b)$.  Intuitively, $a \opG b$ is the grade of
  ``the better of two alternatives.''

\item $\otG : [0,1] \times [0,1] \to [0,1]$ is the
  \textbf{composition} (or \textbf{conjunction}) operation:
  $a \otG b = \min(a, b)$.  Intuitively, $a \otG b$ is the grade of
  ``both $a$ and $b$ together''---the bottleneck principle, where the
  weakest link determines the overall quality.

\item $\gzero = 0$ is the \textbf{zero} element: the grade of an
  artifact that is completely unsatisfactory.

\item $\gone = 1$ is the \textbf{unit} element: the grade of a
  perfect artifact.
\end{enumerate}
\end{definition}

\begin{proposition}[Semiring axioms]
\label{prop:semiring-axioms}
$(\Grade, \opG, \otG, \gzero, \gone)$ satisfies the semiring axioms:
\begin{enumerate}
\item $(\Grade, \opG, \gzero)$ is a commutative monoid:
  \begin{itemize}[nosep]
  \item Associativity: $(a \opG b) \opG c = a \opG (b \opG c)$.
  \item Commutativity: $a \opG b = b \opG a$.
  \item Identity: $a \opG \gzero = a$.
  \end{itemize}

\item $(\Grade, \otG, \gone)$ is a commutative monoid:
  \begin{itemize}[nosep]
  \item Associativity: $(a \otG b) \otG c = a \otG (b \otG c)$.
  \item Commutativity: $a \otG b = b \otG a$.
  \item Identity: $a \otG \gone = a$.
  \end{itemize}

\item Distributivity: $a \otG (b \opG c) = (a \otG b) \opG (a \otG c)$.

\item Absorption: $a \otG \gzero = \gzero$.
\end{enumerate}
\end{proposition}

\begin{proof}
Each axiom follows from the properties of $\max$ and $\min$ on $[0,1]$.
\begin{enumerate}
\item $\max$ is associative, commutative, and $\max(a, 0) = a$ for
  $a \in [0,1]$.

\item $\min$ is associative, commutative, and $\min(a, 1) = a$ for
  $a \in [0,1]$.

\item $\min(a, \max(b, c)) = \max(\min(a, b), \min(a, c))$.
  This is the distributive law for lattices, which holds in any
  distributive lattice.

\item $\min(a, 0) = 0$.  \qedhere
\end{enumerate}
\end{proof}


%% ─── Section 2: Partial order ─────────────────────────────────────

\section{The partial order and lattice structure}
\label{sec:grade:order}

\begin{proposition}[Induced partial order]
\label{prop:grade-order}
The semiring $\Grade$ carries a natural partial order defined by:
\[
  a \leq b \quad \Longleftrightarrow \quad a \opG b = b.
\]
This coincides with the usual order on $[0, 1]$: $a \leq b$ iff
$\max(a, b) = b$ iff $a \leq b$ in $\mathbb{R}$.
\end{proposition}

\begin{proof}
$a \opG b = \max(a, b) = b$ iff $a \leq b$ in the usual order on
$\mathbb{R}$.
\end{proof}

\begin{proposition}[Bounded distributive lattice]
\label{prop:grade-lattice}
Under the order $\leq$, $\Grade$ is a bounded distributive lattice
with:
\begin{itemize}[nosep]
\item Join: $a \vee b = a \opG b = \max(a, b)$.
\item Meet: $a \wedge b = a \otG b = \min(a, b)$.
\item Bottom: $\bot = \gzero = 0$.
\item Top: $\top = \gone = 1$.
\end{itemize}
Distributivity ($a \wedge (b \vee c) = (a \wedge b) \vee (a \wedge c)$)
holds, as verified in \Cref{prop:semiring-axioms}.
\end{proposition}


%% ─── Section 3: Quantale structure ────────────────────────────────

\section{Quantale structure}
\label{sec:grade:quantale}

The Grade semiring is not merely a lattice; it has the richer structure
of a \emph{quantale}: a complete lattice equipped with an associative
multiplication that distributes over arbitrary joins.

\begin{definition}[Quantale]
\label{def:quantale}
A \emph{quantale} is a triple $(Q, \bigvee, \otimes)$ where:
\begin{enumerate}[label=(\roman*),nosep]
\item $(Q, \leq)$ is a complete lattice (every subset has a supremum).
\item $\otimes : Q \times Q \to Q$ is an associative binary operation.
\item $\otimes$ distributes over arbitrary suprema:
  $a \otimes \bigvee S = \bigvee \{a \otimes s \mid s \in S\}$
  and
  $(\bigvee S) \otimes a = \bigvee \{s \otimes a \mid s \in S\}$.
\end{enumerate}
A quantale is \emph{commutative} if $\otimes$ is commutative, and
\emph{unital} if there exists $e \in Q$ with $e \otimes a = a = a \otimes e$
for all $a$.
\end{definition}

\begin{theorem}[$\Grade$ is a commutative unital quantale]
\label{thm:grade-quantale}
$([0,1], \bigvee, \otG)$ is a commutative unital quantale with
unit $\gone = 1$.
\end{theorem}

\begin{proof}
The interval $[0, 1]$ is a complete lattice under the usual order,
with $\bigvee S = \sup S$ for any $S \subseteq [0, 1]$.  The
operation $\otG = \min$ is associative and commutative with unit $1$.
Distributivity of $\min$ over arbitrary suprema: for any $a \in [0,1]$
and $S \subseteq [0,1]$,
\[
  \min(a, \sup S)
  = \sup\{\min(a, s) \mid s \in S\}.
\]
This holds because $\min(a, -)$ preserves directed suprema (it is
Scott-continuous) and, more strongly, preserves all suprema since $[0,1]$
is a complete Heyting algebra.
\end{proof}


%% ─── Section 4: Residuation ──────────────────────────────────────

\section{Residuation: grade implication}
\label{sec:grade:residuation}

The quantale structure on $\Grade$ provides a natural notion of
\emph{implication} via residuation.

\begin{definition}[Residuation]
\label{def:residuation}
The \emph{residual} (or \emph{grade implication}) is defined by:
\[
  a \multimap b
  = \bigvee \{c \in [0,1] \mid a \otG c \leq b\}
  = \bigvee \{c \mid \min(a, c) \leq b\}.
\]
Intuitively, $a \multimap b$ is the \emph{best grade that, combined
with $a$, still does not exceed $b$}.
\end{definition}

\begin{proposition}[Computation of residuation]
\label{prop:residuation-compute}
For $a, b \in [0, 1]$:
\[
  a \multimap b =
  \begin{cases}
    1 & \text{if } a \leq b, \\
    b & \text{if } a > b.
  \end{cases}
\]
\end{proposition}

\begin{proof}
If $a \leq b$: then $\min(a, c) \leq a \leq b$ for all $c$, so
$\{c \mid \min(a, c) \leq b\} = [0, 1]$, and $\bigvee [0,1] = 1$.

If $a > b$: then $\min(a, c) \leq b$ iff $c \leq b$ (since
$a > b$, the minimum is determined by $c$ when $c \leq b$, and equals
$a > b$ when $c > a$, and equals $c$ when $b < c \leq a$).  More
precisely: if $c \leq b$, then $\min(a, c) = c \leq b$; if $c > b$,
then $\min(a, c) \geq \min(a, c) > b$ when $c > b$ and $c \leq a$ gives
$\min(a,c) = c > b$.  So $\{c \mid \min(a, c) \leq b\} = [0, b]$, and
$\bigvee [0, b] = b$.
\end{proof}

\begin{proposition}[Adjunction property]
\label{prop:residuation-adjunction}
Residuation satisfies the adjunction:
\[
  a \otG c \leq b \quad \Longleftrightarrow \quad c \leq (a \multimap b).
\]
This makes $\Grade$ a \emph{closed monoidal category} (more precisely,
a closed symmetric monoidal poset).
\end{proposition}

\begin{proof}
$(\Rightarrow)$: If $a \otG c \leq b$, then
$c \in \{c' \mid a \otG c' \leq b\}$, so $c \leq \bigvee\{c' \mid a \otG c' \leq b\} = a \multimap b$.

$(\Leftarrow)$: If $c \leq a \multimap b$, then
$a \otG c \leq a \otG (a \multimap b)$.  We claim
$a \otG (a \multimap b) \leq b$.  If $a \leq b$: $a \otG 1 = a \leq b$.
If $a > b$: $a \otG b = \min(a, b) = b$.  In both cases,
$a \otG (a \multimap b) \leq b$, so $a \otG c \leq b$.
\end{proof}


%% ─── Section 5: The evaluation functor ───────────────────────────

\section{The evaluation functor}
\label{sec:grade:evaluation}

The Grade semiring acquires its semantic content through the
\emph{evaluation functor}, which maps derivations (proofs, evidence
chains) to grades.

\begin{definition}[Evaluation functor]
\label{def:evaluation-functor}
Let $\mathbf{Prf}$ be the category whose objects are evidence items
(proofs, solver certificates, runtime traces, oracle assertions) and
whose morphisms are evidence transformations (composition, weakening,
restriction).  The \emph{evaluation functor} is a functor
\[
  \mu : \mathbf{Prf} \to \Grade
\]
(viewing $\Grade$ as a category via its partial order) that assigns to
each evidence item a grade and preserves the ordering: stronger
evidence receives a higher grade.
\end{definition}

\begin{remark}[Channel-dependent evaluation]
\label{rem:channel-evaluation}
In practice, $\mu$ is channel-dependent:
\begin{itemize}[nosep]
\item Formal proofs: $\mu(\text{proof}) = \gone = 1$ (maximal confidence).
\item Solver certificates: $\mu(\text{solver}) \in [0.9, 1]$
  (high confidence, modulated by fragment coverage).
\item Runtime witnesses: $\mu(\text{runtime}) \in [0.5, 0.9]$
  (confidence depends on test coverage).
\item Oracle assertions: $\mu(\text{oracle}) \in [0, 0.5]$
  (low baseline confidence, subject to promotion with justification).
\end{itemize}
These ranges are illustrative; the specific evaluation function is a
parameter of the framework, not a fixed choice.
\end{remark}


%% ─── Section 6: Harmony ──────────────────────────────────────────

\section{Harmony: the bottleneck principle}
\label{sec:grade:harmony}

A central application of the Grade semiring is the \emph{harmony}
functional, which computes the overall quality of an artifact by
composing grades across strata.

\begin{definition}[Harmony]
\label{def:harmony}
Let $\varphi$ be an artifact analyzed across $n$ strata
$s_1, \ldots, s_n$ (e.g., in a program: type safety, behavioral
correctness, performance, documentation; in a poem: meter, rhyme,
imagery, emotion, unity).  The \emph{harmony} of $\varphi$ is:
\[
  h(\varphi) = \bigotimes_{s \in \text{strata}} T_s(\varphi)
  = T_{s_1}(\varphi) \otG T_{s_2}(\varphi) \otG \cdots \otG T_{s_n}(\varphi)
  = \min(T_{s_1}(\varphi), \ldots, T_{s_n}(\varphi)).
\]
This is the \textbf{bottleneck principle}: the overall quality is
determined by the weakest stratum.
\end{definition}

\begin{proposition}[Properties of harmony]
\label{prop:harmony-properties}
\begin{enumerate}
\item $h(\varphi) = \gone$ if and only if $T_s(\varphi) = \gone$ for
  every stratum $s$.  (An artifact is perfect only if every aspect is
  perfect.)

\item $h(\varphi) = \gzero$ if and only if $T_s(\varphi) = \gzero$
  for some stratum $s$.  (A single catastrophic failure in any
  stratum renders the artifact worthless.)

\item $h$ is monotone in each argument: improving any stratum's grade
  cannot decrease the harmony.

\item $h$ is idempotent on uniform grades: if $T_s(\varphi) = t$ for
  all $s$, then $h(\varphi) = t$.
\end{enumerate}
\end{proposition}

\begin{proof}
All four properties follow from the corresponding properties of $\min$
on $[0, 1]$.
\end{proof}


%% ─── Section 7: OT as degenerate case ────────────────────────────

\section{Optimality Theory as a degenerate case}
\label{sec:grade:ot}

Optimality Theory (OT), the dominant framework in generative phonology,
can be recovered as a degenerate case of the Grade semiring.

\begin{proposition}[OT as a special case of $\Grade$]
\label{prop:ot-special-case}
In classical OT, constraints are ranked $C_1 \gg C_2 \gg \cdots \gg C_n$,
and each constraint assigns a binary pass/fail to a candidate.  This is
the special case of $\Grade$ where:
\begin{enumerate}[nosep]
\item The grade set is $\{0, 1\}$ (binary).
\item The strata are the OT constraints $C_1, \ldots, C_n$.
\item The harmony functional $h$ evaluates candidates by the
  lexicographic order induced by the constraint ranking.
\end{enumerate}
\end{proposition}

\begin{proof}[Proof sketch]
In binary $\Grade$, $\otG = \min$ on $\{0, 1\}$ is conjunction
(both must pass) and $\opG = \max$ on $\{0, 1\}$ is disjunction
(at least one must pass).  The OT evaluation of a candidate is
the lexicographically ordered tuple $(C_1(\varphi), \ldots, C_n(\varphi))$,
which corresponds to the graded harmony $h(\varphi) =
C_1(\varphi) \otG \cdots \otG C_n(\varphi)$ in the binary Grade
with the lexicographic tiebreaking convention.

The key difference is that classical OT is restricted to binary grades
and strict ranking, while $\Grade$ allows continuous grades and the
quantale structure provides residuation (``how much violation of
$C_i$ can be tolerated given the grade on $C_j$?''), which OT
cannot express.
\end{proof}

\begin{remark}
\label{rem:ot-generalization}
This embedding of OT into $\Grade$ is significant for linguistic
applications of Judgment Geometry.  It means that the JuGeo framework
can express everything OT can express (ranked constraints with binary
evaluations) while also expressing graded constraints, continuous
quality measures, and the residuation calculus.  The multi-stratum
architecture of the Compositional Theory of Theories (Part~III)
exploits this generalization.
\end{remark}


%% ─── Section 8: The trust algebra ────────────────────────────────

\section{The trust algebra: grades along chains}
\label{sec:grade:trust}

When judgments are composed along chains of evidence---one source
produces evidence that is consumed by another---the grades compose
multiplicatively.

\begin{definition}[Trust composition]
\label{def:trust-composition}
If $A$ and $B$ are judgments with grades $T(A)$ and $T(B)$, and $B$
depends on $A$ (i.e., $B$ uses $A$ as evidence), then the grade of
the composite judgment is:
\[
  T(A \circ B) = T(A) \otG T(B) = \min(T(A), T(B)).
\]
This is the \textbf{conservative principle}: a chain of evidence is
only as strong as its weakest link.
\end{definition}

\begin{proposition}[Properties of trust composition]
\label{prop:trust-properties}
\begin{enumerate}
\item \textbf{Monotonicity:} strengthening any link in the chain
  cannot weaken the composite:
  $T(A) \leq T(A')$ implies $T(A \circ B) \leq T(A' \circ B)$.

\item \textbf{Conservative join:} composing a high-grade judgment
  with a low-grade judgment yields the low grade:
  $\gone \otG t = t$ for any $t$.

\item \textbf{Associativity:} $T((A \circ B) \circ C) = T(A \circ (B \circ C))$.

\item \textbf{No silent promotion:} there is no operation that
  increases the grade of a composite beyond the minimum of its
  components, unless an explicit promotion step is recorded in the
  proof trail $\Pi$.  Any system claiming to promote trust must
  justify the promotion and record it.
\end{enumerate}
\end{proposition}

\begin{proof}
Properties (1)--(3) follow from the corresponding properties of $\min$
and the associativity of $\otG$.  Property (4) is a design invariant
of the framework: promotion is a distinguished operation, not a
consequence of the algebra, and its occurrence is always logged in the
proof trail.
\end{proof}

\begin{remark}[The trust algebra as enrichment]
\label{rem:trust-enrichment}
The trust composition law makes the judgment category
$\mathbf{Jud}(\mathcal{C}, J, \mathcal{F})$ into a category
\emph{enriched} over $\Grade$: the hom-objects are not merely sets but
carry grades, and composition of morphisms respects the Grade
multiplication.  This enrichment is the bridge to the Graded
Evaluation object $\GEval$ of Part~II.
\end{remark}


%% ─── Section 9: Examples ──────────────────────────────────────────

\section{Examples of graded evaluation}
\label{sec:grade:examples}

We close with concrete examples of graded evaluation across the four
domains.

\begin{example}[Programs]
\label{ex:grade:programs}
A program is evaluated across three strata:
\begin{itemize}[nosep]
\item $T_{\text{type}} = 1.0$: the program type-checks (formal proof).
\item $T_{\text{behav}} = 0.85$: the program passes 85\% of behavioral
  test cases (runtime evidence).
\item $T_{\text{perf}} = 0.72$: the program meets 72\% of performance
  benchmarks (runtime evidence).
\end{itemize}
The harmony is $h = \min(1.0, 0.85, 0.72) = 0.72$.  The bottleneck
is performance; improving performance would have the greatest impact
on overall quality.  The residual $T_{\text{perf}} \multimap T_{\text{type}} = 1.0$
tells us that performance quality does not constrain type safety---they are
independent strata.
\end{example}

\begin{example}[Natural language]
\label{ex:grade:nl}
A sentence is evaluated across four strata:
\begin{itemize}[nosep]
\item $T_{\phon} = 0.95$: phonological well-formedness.
\item $T_{\synt} = 1.0$: syntactic grammaticality.
\item $T_{\sem} = 0.80$: semantic coherence.
\item $T_{\prag} = 0.60$: pragmatic felicity.
\end{itemize}
The harmony is $h = \min(0.95, 1.0, 0.80, 0.60) = 0.60$.  The
bottleneck is pragmatic felicity: the sentence is grammatical and
meaningful but contextually odd.  The residual
$T_{\prag} \multimap T_{\synt} = 1.0$ confirms that pragmatic
infelicity does not impugn syntactic grammaticality.
\end{example}

\begin{example}[Poetry]
\label{ex:grade:poetry}
A sonnet is evaluated across five strata:
\begin{itemize}[nosep]
\item $T_{\text{meter}} = 0.90$: metrical regularity.
\item $T_{\text{rhyme}} = 0.95$: rhyme scheme adherence.
\item $T_{\text{imagery}} = 0.85$: vividness and originality of imagery.
\item $T_{\text{emotion}} = 0.70$: emotional resonance.
\item $T_{\text{unity}} = 0.55$: thematic unity across stanzas.
\end{itemize}
The harmony is $h = \min(0.90, 0.95, 0.85, 0.70, 0.55) = 0.55$.
The bottleneck is thematic unity: the individual stanzas are
accomplished but the poem lacks a coherent arc.  This is precisely
the kind of defect detected by the \v{C}ech cohomology of
\Cref{ch:cech}---the local sections (stanza-level quality) are
high, but they fail to glue (the volta fails, or the thematic
thread is broken).
\end{example}

\begin{example}[Mathematics]
\label{ex:grade:math}
A mathematical proof is evaluated across four strata:
\begin{itemize}[nosep]
\item $T_{\text{rigor}} = 1.0$: every step is formally justified.
\item $T_{\text{elegance}} = 0.65$: the proof is correct but
  ungainly (case-bashing rather than conceptual argument).
\item $T_{\text{novelty}} = 0.80$: the result extends known work in
  a non-trivial way.
\item $T_{\text{significance}} = 0.90$: the result has important
  consequences.
\end{itemize}
The harmony is $h = \min(1.0, 0.65, 0.80, 0.90) = 0.65$.  The
bottleneck is elegance.  The proof is unimpeachably correct but fails
to achieve the aesthetic standard expected of a publishable result.
The residual $T_{\text{elegance}} \multimap T_{\text{rigor}} = 1.0$
confirms that lack of elegance does not compromise rigor.
\end{example}

\begin{remark}[The interaction of grade and cohomology]
\label{rem:grade-cohomology-interaction}
The Grade semiring and the \v{C}ech cohomology interact in a
fundamental way.  When $\check{H}^1 \neq 0$, the obstruction class
$B$ depresses the grade: the overall harmony cannot exceed the grade
of the worst overlap.  Conversely, when $\check{H}^1 = 0$ and all
local grades are high, the global grade is high.  This interaction
is formalized in Part~II via the Graded Evaluation object $\GEval$,
which pairs the quantale structure of $\Grade$ with the spectral
sequence of $\CSpec$ to produce a single coherent quality assessment.
\end{remark}
